% concatenated from D260215.tex
\documentclass[11pt]{article}
\usepackage[ngerman, croatian, english]{babel}
\usepackage[T1]{fontenc}
\usepackage[utf8]{inputenc}
\usepackage{geometry}
\geometry{verbose,tmargin=2cm,bmargin=2.5cm,lmargin=3cm,rmargin=3cm}
\usepackage{color}
\usepackage{babel}
\usepackage{textcomp}
\usepackage{mathrsfs}
\usepackage{amsmath}
\usepackage{amsthm}
\usepackage{amssymb}
\usepackage{cancel}
\usepackage{stmaryrd}
\usepackage{stackrel}
\usepackage[all]{xy}
\usepackage[unicode=true,pdfusetitle,
bookmarks=true,bookmarksnumbered=false,bookmarksopen=false,
breaklinks=false,pdfborder={0 0 1},backref=false,colorlinks=true]
{hyperref}
\hypersetup{
	pdfborderstyle=,linkcolor=MidnightBlue, citecolor=PineGreen, urlcolor=PineGreen}

\makeatletter
%%%%%%%%%%%%%%%%%%%%%%%%%%%%%% Textclass specific LaTeX commands.
\theoremstyle{plain}
\newtheorem{thm}{\protect\theoremname}[section]
\theoremstyle{definition}
\newtheorem{defn}[thm]{\protect\definitionname}
\theoremstyle{remark}
\newtheorem{rem}[thm]{\protect\remarkname}
\theoremstyle{plain}
\newtheorem{lem}[thm]{\protect\lemmaname}
\theoremstyle{plain}
\newtheorem{prop}[thm]{\protect\propositionname}
\theoremstyle{plain}
\newtheorem{cor}[thm]{\protect\corollaryname}
\theoremstyle{definition}
\newtheorem*{problem*}{\protect\problemname}
\newtheorem*{funding*}{\protect\fundingname}

\@ifundefined{date}{}{\date{}}
%%%%%%%%%%%%%%%%%%%%%%%%%%%%%% User specified LaTeX commands.
\usepackage[all]{xy}
\usepackage[dvipsnames]{xcolor}
\usepackage{tikz}
\usepackage{hyperref}

\usepackage{csquotes}
\usepackage[backend=biber, style=alphabetic, maxnames=5]{biblatex}
\DeclareFieldFormat{urldate}{}
\DeclareFieldFormat{language}{}
\DeclareFieldFormat{url}{}
\addbibresource{Irreducible.bib}


\usepackage{changepage}

\usepackage{graphicx}
\usepackage{xstring}
\usepackage{forloop}
\usepackage{bbm, amsmath, amsfonts, amscd, latexsym, amsthm, amssymb, graphicx, mathrsfs, comment}
\usepackage{mathtools}
\usepackage{caption}
\usepackage{MnSymbol}
\usepackage{booktabs}


\renewcommand\texteuro{FIXME}

\makeatother

\providecommand{\corollaryname}{Corollary}
\providecommand{\definitionname}{Definition}
\providecommand{\lemmaname}{Lemma}
\providecommand{\problemname}{Problem}
\providecommand{\propositionname}{Proposition}
\providecommand{\remarkname}{Remark}
\providecommand{\theoremname}{Theorem}

\begin{document}
\title{Components of discriminants for systems of equations and irreducibility
of determinants}
\author{Vladislav Pokidkin\thanks{Email: vppokidkin@hse.ru\protect \\
National Research University Higher School of Economics, Russian Federation\protect \\
AG Laboratory, HSE, 6 Usacheva str., Moscow, Russia, 119048}}
\maketitle
\begin{abstract}
The discriminant of a multivariate polynomial with indeterminate coefficients
is not necessarily a hypersurface, and characterizing its codimension
was an open problem for quite a while. We resolve this problem for
the discriminants of systems of polynomials with indeterminate coefficients
and with the same number of equations and unknowns (square polynomial
systems). This version is more involved in the sense that the discriminant
may have several components of different dimensions. 

In the space of square matrices, we characterize row-generated subspaces
on which the determinant is an irreducible polynomial. This allows
us to resolve the Esterov conjecture for square polynomial systems
whose discriminant is an irreducible hypersurface. Based on this result,
we enumerate all the components and determine their dimensions and
degrees for each of the three conventional ways to formalize the notion
of a discriminant in this setting (mixed, Cayley, and A-discriminants)
in cases of square and overdetermined systems. The proof of Esterov's
conjecture and descriptions of the three types of discriminants are
based on the theory of polymatroids.
\end{abstract}
\textbf{Keywords:} determinants, Esterov conjecture, discriminants
of polynomial systems, Cayley discriminants, mixed discriminants,
realizable polymatroids.

\section{Introduction}

Given an algebraic torus $T\simeq(\mathbb{C}^{\times})^{n}$ with
its character lattice $M\simeq\mathbb{Z}^{n}$, a finite set of monomials
$A\subset M$ generates a vector space of Laurent polynomials, denoted
by $\mathbb{C}_{A}$. For a tuple of finite sets $\mathscr{A}=(A_{1},\ldots,A_{k})\subset M$,
starting from \cite{gelfand_discriminants_1994}, lots of attention
is paid to the $\mathscr{A}$\textit{-discriminant} $D_{\mathscr{A}}\subset\mathbb{C}_{\mathscr{A}}=\mathbb{C}_{A_{1}}\oplus\cdots\oplus\mathbb{C}_{A_{k}}$,
the closure of all tuples of polynomials $(f_{1},\ldots,f_{k})$ such
that the system of equations $f_{1}=\cdots=f_{k}=0$ has a \textit{degenerate
root} (i.e. a point $x\in T$ at which $f_{1}(x)=\cdots=f_{k}(x)=0$
and $df_{1}\wedge\ldots\wedge df_{k}(x)=0$). The motivation varies
from algebraic geometry, algebraic statistics \cite{huh_likelihood_2014,amendola_maximum_2019}
and PDEs \cite{gelfand_hypergeometric_1989,gelfand_generalized_1990}
to mathematical physics \cite{vanhove_feynman_2019,mizera_landau_2022,matsubara-heo_four_2023}
and symbolic algebra \cite{sturmfels_solving_2002}.

A question of particular interest is the classification of tuples
$\mathscr{A}$, for which the discriminant is not a hypersurface.
For instance, the case $k=1$ is equivalent to the classical problem
of dual defective varieties (whose projectively dual is not a hypersurface),
for the special case of toric varieties. This problem was resolved
in \cite{di_rocco_projective_2006,curran_restriction_2007,dickenstein_tropical_2007,esterov_newton_2010,matsui_geometric_2011,furukawa_combinatorial_2021,cattani_non-splitting_2022}.
At the other extreme, for $n=k$, the discriminant was shown to have
a hypersurface component if the support sets $A_{1},\ldots,A_{n}$
cannot be shifted to an affine plane or the tuple of standard simplexes
by an automorphism of the lattice. This was proved in \cite{esterov_galois_2019,borger_defectivity_2020},
motivated by applications in Galois theory and lattice polytope geometry
respectively.

It is possible to define two more types of discriminants. For a tuple
$\mathscr{A},$ the \textit{mixed discriminant} is the closure of
all polynomial systems in $\mathbb{C}_{\mathscr{A}}$ having a \textit{non-degenerate
multiple root} (i.e. a degenerate root $x\in T$ such that no proper
subtuple of $df_{1},...,df_{k}(x)$ is linearly dependent). Mixed
discriminants were introduced in \cite{cattani_mixed_2013} and were
investigated in \cite{dokken_plane_2014,esterov_galois_2019,dickenstein_iterated_2023}.

For a subset $I\subseteq\{1,\ldots,k\}$, the \textit{Cayley trick}
for the subtuple $\mathscr{B}=(A_{i},\;i\in I)$ is a map sending
a polynomial system $f\in\mathbb{C}_{\mathscr{A}}$ to the polynomial
$\sum_{i\in I}\lambda_{i}f_{i}(x)$ in variables $x$ and $\lambda$.
The support of this polynomial is the \textit{Cayley set} $cay(\mathscr{B})=\cup_{i\in I}A_{i}\times\{i\}\subset M\times\mathbb{Z}^{|I|}$.
The \textit{Cayley discriminant} of a subtuple $\mathscr{B}$ is the
preimage of the discriminant for the Cayley set $cay(\mathscr{B})$
under the Cayley trick \cite{esterov_newton_2010,esterov_galois_2019}. The Cayley discriminant was introduced as an intermediate object allowing one to reduce the study of $\mathscr{A}$-discriminants of polynomial
systems to $A$-discriminants of one polynomial (which is much simpler
and well-understood).

In case $n=k$, if the Cayley discriminant is a hypersurface, then
the mixed discriminant is the same hypersurface \cite{cattani_mixed_2013}.
For tuples, called irreducible, Esterov showed that the three types
of discriminants have the same hypersurface component and conjectured
a lack of other components for $\mathscr{A}$-discriminants \cite{esterov_galois_2019}.
We prove the conjecture, and the three types of discriminants are
the same hypersurface for an irreducible tuple.

The following result of independent interest was crucial for the proof
of Esterov's irreducibility theorem. In an $n$-dimensional vector
space $V$ over an algebraically closed field with the Zariski topology,
a tuple of vector subspaces $(L_{1},...,L_{n})$ is \textit{irreducible}
if no $k$ of them lie in the same $k$-dimensional subspace for $0<k<n$.
Then the product of vector spaces $V^{\times n}$ is isomorphic to
the space of $n\times n$ square matrices. The space of square matrices
contains the determinant hypersurface --- the set of matrices with
a zero determinant.
\begin{thm}
For an irreducible tuple $(L_{1},...,L_{n})$ of vector subspaces
in an $n$-dimensional vector space over an algebraically closed field,
the intersection of the determinant hypersurface with the subspace
$L_{1}\times...\times L_{n}$ is irreducible in the space of $n\times n$
square matrices \textup{(Theorem \ref{Theorem. Intersection of the determinant with a subspace})}.
\end{thm}

The proof is based on the polymatroid partition of the dual vector
space \cite{pokidkin_combinatorics_2025}. 

When $1<k<n$, the study of discriminants is substantially more complicated
than in the classical case $k=1$ and in our case $n\leq k$: see
\cite{esterov_newton_2010,esterov_discriminant_2013,dickenstein_iterated_2023}
for some partial results.

We study the case $n\leq k$ more comprehensively: for arbitrary support
sets $\mathscr{A}=(A_{1},\ldots,A_{k})\subset M$, we give the complete
list of the irreducible components of the $\mathscr{A}$-discriminant,
the Cayley discriminant, the mixed discriminant in $\mathbb{C}_{\mathscr{A}}$,
specifying their codimensions and degrees. It turns out that the realizable
polymatroid associated with the tuple $\mathscr{A}$ encodes each
of the three types of discriminants and the resultant. This link between
the theory of polymatroids and algebraic geometry differs from the
one established in works \cite{pagaria_hodge_2023,crowley_bergman_2024}. 

The answer is stated in terms of the following fundamental quantity:
the \textit{defect} of a subtuple $\mathscr{B}$ is the number $\delta(\mathscr{B})=\dim($the
affine span of the Minkowski sum $\sum_{i\in I}A_{i})-|I|$. A tuple
is\textit{ dependent} if it contains a subtuple with negative defect.

We first describe discriminants for dependent tuples: this simplest
case reduces to the sparse resultant as introduced in \cite{sturmfels_newton_1994}. 
\begin{thm}
\label{Theorem. Introduction. Discriminants for linearly dependent tuples}For
a dependent tuple $\mathscr{A}$ with the subtuple $\mathscr{M}$
that has minimal defect and is minimal by inclusion \textup{(Theorem
\ref{Theorem. A/M  is linearly independent})}, 

- the $\mathscr{A}$-discriminant is the sparse resultant $R_{\mathscr{M}}$
of codimension $-\delta(\mathscr{M})$\textup{ (Corollary \ref{Corollary. Discriminants for linearly dependent tuples})}; 

- the mixed discriminant is empty if the defect of the subtuple $\mathscr{M}$
is less than $-1$; otherwise, the mixed discriminant is the sparse
resultant $R_{\mathscr{M}}$, and it is a hypersurface\textup{ (Theorem
\ref{Theorem. The mixed discriminant for a linearly dependent tuple}).}
\end{thm}

We provide a new characterization of the subtuple $\mathscr{M}$ for
the resultant as the maximal cycle of the induced matroid from the
realizable polymatroid on $\mathscr{A}$ (Theorem \ref{Theorem. A/M  is linearly independent}).
We also link mixed discriminants with circuits of the induced matroid
for a dependent tuple $\mathscr{A}$.

The case of independent tuples is the essence of the matter, and we
start with irreducible tuples. A tuple of finite sets is \textit{irreducible}
if the defects of all proper subtuples are positive. Independent tuples
of zero defect are called \textit{BK-tuples}. We call a tuple \textit{linear}
if it can be mapped to the tuple of standard simplexes by an automorphism
of the lattice. In 2018, Esterov conjectured \cite{esterov_galois_2019}:
\begin{thm}
\label{Theorem. Esterov conjecture}For a nonlinear irreducible BK-tuple
$\mathscr{A}$, the $\mathscr{A}$-discriminant is a hypersurface
in the space of polynomial systems $\mathbb{C}_{\mathscr{A}}$ \textup{(Theorem
\ref{Theorem. Codimension of discriminants})}.
\end{thm}

For an irreducible BK-tuple $\mathscr{A}$, the monodromy group is
the symmetric group on the mixed volume $\mathrm{MV}(\mathscr{A})$
elements \cite{esterov_galois_2019}, and the general polynomial system
is solvable by radicals if the mixed volume $\mathrm{MV}(\mathscr{A})$
does not exceed four \cite{esterov_multivariate_2016}. However, these
questions remain open for reducible BK-tuples, and the current work
facilitates their solution in future research.

For a reducible BK-tuple $\mathscr{A}$, Esterov computed the hypersurface
components of the $\mathscr{A}$-discriminant:
\begin{thm}
\textup{(Esterov)\label{Theorem. Esterov, hypersurface components}}
For a reducible BK-tuple $\mathscr{A}$, the codimension one components
of the $\mathscr{A}$-discriminant are Cayley discriminants of BK-subtuples
\textup{(Theorem 2.31, \cite{esterov_newton_2010})}.
\end{thm}

In the current work, we specify Theorem \ref{Theorem. Esterov, hypersurface components}.
The general description of discriminants for square polynomial systems
is based on Theorem \ref{Theorem. Esterov conjecture} and requires
the following notions.
\begin{defn}
A BK-tuple is \textit{simple} if it is not a union of its proper BK-subtuples.
A simple BK-subtuple is \textit{maximal} if it is not contained in
another simple BK-subtuple.

A simple BK-tuple is\textit{ prelinear} if, for every projection of
lattices, sending every set from its maximal (by inclusion) proper
BK-subtuple to zero, the image is a linear tuple.
\end{defn}

\begin{thm}
\label{Theorem. Introduction. Discriminants for linearly independent tuples}
For a BK-tuple $\mathscr{A}$, the $\mathscr{A}$-discriminant has
the following distinct components:

- codimension one Cayley discriminants of non-prelinear simple BK-subtuples;

- codimension two Cayley discriminants of prelinear simple BK-subtuples
not contained in non-prelinear BK-subtuples \textup{(Theorem \ref{Theorem. Discriminant for a BK-tuple}
and Proposition \ref{Proposition. Discriminant's strata are Cayley discriminants})}.
\end{thm}

\begin{rem}
1) If a prelinear BK-subtuple is contained in a non-prelinear BK-subtuple,
then the Cayley discriminant of the former subtuple lies in the Cayley
discriminant of the latter.

2) The proof is based on a specific partition of the tuple $\mathscr{A}$
from Corollary \ref{Corollary. Unique decomposition of a BK-tuple}.
This partition originates from a partition of the realizable polymatroid
on $\mathscr{A}$ \cite{pokidkin_combinatorics_2025}.
\end{rem}

\begin{thm}
For a BK-tuple, the mixed discriminant is empty/a hypersurface/a variety
of codimension two if the tuple is non-simple/non-prelinear simple/prelinear
simple. For a simple BK-tuple, the mixed discriminant equals the Cayley
discriminant \textup{(Theorem \ref{Theorem. The mixed discriminant for a BK-tuple};
\cite{cattani_mixed_2013})}.
\end{thm}

The known results \cite{pedersen_product_1993,gelfand_discriminants_1994,dickenstein_tropical_2007,esterov_determinantal_2007,matsui_geometric_2011}
allow us to write degrees for discriminants. For a dependent tuple/BK-tuple
$\mathscr{A}$, see Propositions \ref{Proposition. I Degree of the sparse resultant for a linearly dependent tuple},
\ref{Proposition. II Degree of the sparse resultant for a linearly dependent tuple}/Corollary
\ref{Corollary. Degrees for components of the discriminant for a BK-tuple}
and Corollary \ref{Corollary. Degree of the mixed discriminant for a linearly dependent tuple}/Remark
\ref{Remark. Degree of the mixed discriminant for a BK-tuple} for
degrees of the $\mathscr{A}$-discriminant and the mixed discriminant
respectively.

Despite Cayley discriminants being well-known $A$-discriminants,
the current paper provides an alternative view of Cayley discriminants
for  dependent and BK-tuples in Theorems \ref{Theorem. The Cayley discriminant for a linearly dependent tuple}
and \ref{Theorem. Cayley Discriminants}. For a BK-tuple, Theorem
\ref{Theorem. Cayley Discriminants} shows the Cayley discriminant
equals the complete intersection of Cayley discriminants of all maximal
simple BK-subtuples. This theorem leads to a slight simplification
of the Matsui-Takeuchi degree formula in Corollary \ref{Corollary. Degree of the Cayley discriminant for a BK-tuple}.

Defining equations of discriminants can be found using computer algebra
systems and software (\href{https://symbolics.juliasymbolics.org/dev}{Julia},
\href{https://macaulay2.com}{Macaulay2}, \href{https://www.oscar-system.org}{Oscar},
\href{https://www.sagemath.org}{Sage}, ...), and then used to enumerate
the components of the discriminants and determine their geometric
characteristics. However, the complexity of this computation rapidly
grows with dimension and the size of the support sets. Our results
express the geometry of the discriminant components combinatorially
in terms of the support sets, circumventing the costly symbolic computation
of discriminants. 

The structure of the paper is as follows. Section \ref{Section. Irreducible intersections of subspaces with determinant}
characterizes irreducible intersections for row-generated subspaces
with the determinant hypersurface in the space of square matrices.
Section \ref{Section. Esterov irreducibility theorem} is dedicated
to the Esterov conjecture. Sections \ref{Section. Combinatorial Review}
and \ref{Section. Discriminants-of-polynomial} remind us some combinatorial
results and general facts about discriminants. In Section \ref{Section. BK-multiplicaiton},
we construct a special multiplication of varieties necessary to describe
the discriminants for BK-tuples. Section \ref{Section. Discriminants for BK-tuples}
characterizes $\mathscr{A}$-discriminants, Section \ref{Section. Cayley discriminants}
- Cayley discriminants, and Section \ref{Section. Mixed discriminants}
- mixed discriminants. For each type of discriminants, we enumerate
components and compute their codimensions first for BK-tuples and
then for dependent tuples. All computations of degrees are collected
in Section \ref{Section. Degrees of components}.

\section{\label{Section. Irreducible intersections of subspaces with determinant}Irreducible
intersections of subspaces with determinant}

For a quasi-affine algebraic set $X$, we denote its Zariski closure
by $\overline{X}$. An irreducible quasi-affine algebraic set over
an algebraically closed field we call a \textit{variety}. A \textit{hypersurface}
is a codimension one variety.

Consider an $n$-dimensional vector space $V$ over an algebraically
closed field $\mathbb{K}$ with the Zariski topology. Let $E$ be
an algebraic set defined by the equation $\varphi(l)=0$ in the space
$\mathrm{End}(V^{\lor})\times V^{\lor}\ni(\varphi,l),$ where $V^{\lor}$
is the dual space, and $\mathrm{End}(V^{\lor})$ is the vector space
of endomorphisms of $V^{\vee}$. By choosing coordinates, the equation
$\varphi(l)=0$ determines an intersection of $n$ quadrics in $\mathrm{End}(V^{\lor})\times V^{\lor}$.
For a nontrivial covector $l$, every point $(\varphi,l)$ of $E$
corresponds to a degenerate linear operator $\varphi$ with an eigenvector
$l$ of eigenvalue zero. 

We can explicitly write the quadric equations. Choose a basis $e_{1},...,e_{n}$
in the vector space $V$, the dual basis $e^{1},...,e^{n}$ in the
dual space $V^{\vee}$, and the basis $e^{i}\otimes e_{j}$ in the
space $\mathrm{End}(V^{\lor})$ using the isomorphism $\mathrm{End}(V^{\lor})\cong V^{\lor}\otimes V$.
Following the Einstein convention, we can represent the elements of
vector spaces in coordinates: $l=l_{i}e^{i}$ and $\varphi=\varphi_{i}^{j}e^{i}\otimes e_{j},$
$l_{i},\varphi_{i}^{j}\in\mathbb{K}$. Then the equation $\varphi(l)=0$
is equivalent to the system $\varphi_{i}^{j}l_{j}e^{i}=0.$ Since
vectors $e^{i}$ form a basis, we need the coefficients $\varphi_{i}^{j}l_{j}$
to be zero. We obtain an intersection of $n$ quadrics $\varphi_{i}^{j}l_{j}=0$
in $\mathrm{End}(V^{\lor})\times V^{\lor}$.

Notice that the equations $\varphi_{i}^{j}l_{j}=0$ are equivalent
to $l(x_{i})=0,$ $x_{i}=\varphi_{i}^{j}e_{j}\in V$. Hence there
exists an isomorphism $\mathrm{End}(V^{\lor})\cong V^{\times n}$
such that $\varphi(l)=l(x_{i})e^{i}$ for a tuple of points $(x_{1},...,x_{n})\in V\times...\times V$.
Then the equations $l(x_{i})=0$ define the necessary intersection
$E$ of $n$ quadrics in the space $V^{\times n}\times V^{\vee}.$
Remarkably, we do not need to choose bases in $V$ and $V^{\vee}$
to write the system of equations $l(x_{i})=0,$ $i\in[n]$ in $V^{\times n}\times V^{\vee}\ni(x_{1},...,x_{n},l).$

There are two natural projections: 
\[
\begin{array}[t]{cc}
\xymatrix{ & E\ar@{->>}[ld]_{p}\ar@{->>}[rd]^{q}\\
V^{\times n} &  & V^{\vee}
}
 & \begin{array}[t]{l}
p(x_{1},...,x_{n},l)=(x_{1},...,x_{n}),\\
q(x_{1},...,x_{n},l)=l.
\end{array}\end{array}
\]
For a non-zero covector $l$, the fiber $q^{-1}(l)$ is a vector subspace
in $V^{\times n}$ of dimension $n^{2}-n$. For a matrix $\varphi$,
the fiber $p^{-1}(\varphi)$ is a subspace in $V^{\lor}$ such that
its dimension equals the geometric multiplicity of zero (as an eigenvalue).

Recall that a \textit{polymatroid} is a pair $([n],rk)$ of a finite
set $[n]$ and a rank function $rk:2^{[n]}\rightarrow\mathbb{Z}_{\geq0}$,
which is submodular, monotone, and normalized ($rk(\varnothing)=0$).
The pair of a subspace tuple $(L_{1},...,L_{n})$ from the vector
space $V$ and the rank function $rk_{P}(I)=dim\,L_{I}$, $L_{I}=\underset{i\in I}{\sum}L_{i}$,
forms a polymatroid $P$. The same polymatroid $P$ is defined as
the pair of orthogonal subspaces $(L_{1}^{\perp},...,L_{n}^{\perp})$
in the dual space $V^{\vee}$ and the rank function $rk_{P}^{\perp}(I)=codim\,L_{I}^{\perp},$
$L_{I}^{\perp}=\underset{i\in I}{\cap}L_{i}^{\perp}$. 

The\textit{ defect} of a set $I\subseteq[n]$ is the number $\delta(I)=rk_{P}(I)-|I|$.

A \textit{flat} of the polymatroid $P$ is a subset $F\subseteq[n]$
that is maximal among sets of its rank. The set of all flats forms
an order lattice $\mathcal{L}$ by inclusion. The dual space $V^{\vee}$
admits a polymatroid partition into constructible sets $B_{F}$, enumerated
by flats $F$ from the lattice $\mathcal{L}$ by \cite{pokidkin_combinatorics_2025}.
The sets $B_{F}$ are dense in the subspaces $L_{F}^{\perp}$. Denote
by $L$ the product $L_{1}\times...\times L_{n}\subseteq V^{\times n}$
and by $E_{l}|_{L}$ the restriction of the fiber $q^{-1}(l)|_{L}=p(q^{-1}(l))\cap L$
to the subspace $L$ for a covector $l$.
\begin{lem}
\label{Lemma. Dimension of a fiber}For every $l\in B_{F},$ the restriction's
dimension equals $dim\,E_{l}|_{L}=dim\,L-n+|F|.$
\end{lem}

\begin{proof}
Every constructible set is given by $B_{F}=L_{F}^{\perp}\backslash\underset{F'\in\mathcal{L}\backslash(F)}{\cup}L_{F'}^{\perp}$
for a flat $F$. We can rewrite the definition as follows: $B_{F}=\{l\in V^{\vee}\,|\;l\in L_{F}^{\perp},$
and $l\notin L_{j}^{\perp},\;j\in[n]\backslash F\}=\{l\in V^{\vee}\,|\;l^{\perp}\supseteq L_{F},$
and $l^{\perp}\cancel{\supseteq}L_{j},\;j\in[n]\backslash F\}$. This
means that the set $B_{F}$ corresponds to the set of points $l,$
for which the orthogonal complement $l^{\perp}$ contains each subspace
$L_{i}$ for all indices $i$ from the flat $F$, and the complement
$l^{\perp}$ does not contain subspaces $L_{j}$ for indices $j$
from the complement $[n]\backslash F$. Therefore, the dimension for
the restriction of the fiber $E_{l}|_{L}=l^{\perp}\times...\times l^{\perp}\cap L=(l^{\perp}\cap L_{1})\times...\times(l^{\perp}\cap L_{n})$
to the subspace $L$ equals 
\begin{align*}
dim\,E_{l}|_{L} & =\stackrel[i=1]{n}{\sum}dim\,l^{\perp}\cap L_{i}=\stackrel[i=1]{n}{\sum}(dim\,l^{\perp}+dim\,L_{i}-dim\,(l^{\perp}+L_{i}))=\\
 & =(n^{2}-n)+dim\,L-(n^{2}-n+n-|F|)=dim\,L-n+|F|.
\end{align*}
\end{proof}
Denote by $Q_{F}$ the intersection $q^{-1}(B_{F})\cap L\times V^{\vee}$.
Use the polymatroid partition $V^{\vee}=\underset{F\in\mathcal{L}}{\sqcup}B_{F}$
and Lemma \ref{Lemma. Dimension of a fiber} to write the partition
for the intersection:
\[
E\cap(L\times V^{\vee})=\underset{F\in\mathcal{L}}{\sqcup}Q_{F}.
\]

\begin{prop}
\label{Proposition. Vector bundles other B_F}For a proper flat $F$,
the triple $(Q_{F},q,B_{F})$ is a vector bundle of rank $r=dim\,L-n+|F|$,
where $Q_{F}$ is a variety of dimension $dim\,L-\delta(F)$.
\end{prop}

\begin{proof}
By Lemma \ref{Lemma. Dimension of a fiber}, every fiber is a vector
space of dimension $r$. Therefore, the dimension of the preimage
$Q_{F}$ equals 
\[
dim\,Q_{F}=dim\,B_{F}+dim\,fiber=(n-dim\,L_{F})+(dim\,L-n+|F|)=dim\,L-\delta(F).
\]

Let us show that it is possible to cover the base $B_{F}$ by open
charts $Y_{F}$ such that the preimages $q^{-1}(Y_{F})$ are isomorphic
to trivial vector bundles $Y_{F}\times\mathbb{K}^{r}$. Choose arbitrary
vectors $v_{i}\in L_{i}\backslash L_{F}$ for each element $i\in[n]\backslash F$
and define the set $Y_{F}=B_{F}\backslash\underset{i\in[n]\backslash F}{\cup}v_{i}^{\perp}$,
which is open in the vector subspace $L_{F}^{\perp}$. Choose an arbitrary
decomposition of each subspace $L_{i}$ into the direct sum $L_{i}=U_{i}\oplus\langle v_{i}\rangle$
for each element $i\in[n]\backslash F$. Notice that the preimage
$q^{-1}(Y_{F})$ is an open subset of $Q_{F}$ and equals the intersection
$E\cap L\times Y_{F}$. By the construction of $B_{F},$ the preimage
$Q_{F}$ is an intersection of the set $L\times B_{F}$ and $n-|F|$
quadratic equations of the form: $l(x_{i})=0,$ $x_{i}\in L_{i}$,
$i\in[n]\backslash F$, and $l\in B_{F}$. If we restrict the equations
to the set $L\times Y_{F}$ --- recall the decomposition $L_{i}=U_{i}\oplus\langle v_{i}\rangle$,
$x_{i}=u_{i}+\lambda_{i}v_{i}$, $u_{i}\in U_{i},$ $\lambda_{i}\in\mathbb{K}$
--- we obtain
\[
l(u_{i})+\lambda_{i}l(v_{i})=0,\qquad l\in Y_{F}.
\]
Remarkably, the value of $l(\upsilon_{i})$ is never zero for $l\in Y_{F}$
for each $i$. This means we can express the variables $\lambda_{i}$
in terms of these equations. Then the coordinate algebra of the preimage
$q^{-1}(Y_{F})$ is 
\[
\mathbb{K}[q^{-1}(Y_{F})]\cong\mathbb{K}[\underset{i\in F}{\oplus}L_{i}\oplus\underset{j\in[n]\backslash F}{\oplus}U_{j}\times Y_{F}],
\]
isomorphic to the free algebra over the ring $\mathbb{K}[Y_{F}]$.
Therefore, the preimage $q^{-1}(Y_{F})$ is a variety and a trivial
vector bundle over the open chart $Y_{F}$ with the expected fiber.

Notice that the choice of vectors $v_{i}$ and subspaces $U_{i}$
was arbitrary. Let us show that the base $B_{F}$ can be covered by
affine charts of the form $Y_{F}=Y_{F}(v)$, $v=(v_{i})_{i\in[n]\backslash F}\in C_{F}=\underset{i\in[n]\backslash F}{\prod}L_{i}\backslash L_{F}$,
\begin{align*}
\underset{\upsilon\in C_{F}}{\cup}Y_{F}(v) & =\underset{v\in C_{F}}{\cup}(B_{F}\backslash\underset{i\in[n]\backslash F}{\cup}v_{i}^{\perp})=\underset{\upsilon\in C_{F}}{\cup}\underset{i\in[n]\backslash F}{\cap}B_{F}\backslash v_{i}^{\perp}=\underset{i\in[n]\backslash F}{\cap}(\underset{v_{i}\in L_{i}\backslash L_{F}}{\cup}B_{F}\backslash v_{i}^{\perp})\overset{(*)}{=}B_{F},\\
(*) & \underset{v_{i}\in L_{i}\backslash L_{F}}{\cup}B_{F}\backslash v_{i}^{\perp}=B_{F}\backslash\underset{v_{i}\in L_{i}\backslash L_{F}}{\cap}v_{i}^{\perp}=B_{F}\backslash L_{i}^{\perp}=B_{F},
\end{align*}
\end{proof}
\begin{prop}
The intersection of the determinant hypersurface $Det$ with the subspace
$L$ in the space of square matrices equals the union:
\[
Det\cap L=\underset{F\in\mathcal{L}\backslash[n]}{\cup}\overline{p(Q_{F})}.
\]
\end{prop}

\begin{proof}
The single polynomial equation $det\,\varphi|_{L}$ defines the intersection
$Det\cap L$. Hence the intersection is closed, and all components
have codimension one in the subspace $L$.

By the polymatroid partition of the dual vector space \cite{pokidkin_combinatorics_2025}
and Propositions \ref{Proposition. Vector bundles other B_F}, the
intersection $E\cap L\times V^{\lor}=\underset{F\in\mathcal{L}}{\cup}Q_{F}$
can be written as the union of strata $Q_{F}$, enumerated by flats
$F$ of the polymatroid on $(L_{1},...,L_{n})$ with the lattice of
flats $\mathcal{L}$. By definition, the zero vector cannot be an
eigenvector. This means every point $(\varphi,l)$ from $E\backslash q^{-1}(0)$
is a pair consisting of a degenerate matrix $\varphi$ with a non-zero
eigenvector $l$, corresponding to eigenvalue zero. Then the determinant
hypersurface $Det$ equals the closure of the projection $p(E\backslash q^{-1}(0))$.
Notice that $B_{[n]}=L_{[n]}^{\perp}=(L_{[n]})^{\perp}=V^{\perp}=\{0\}$,
$Q_{[n]}=q^{-1}(B_{[n]})\cap L\times V^{\lor}=L\times\{0\}$, and
the intersection $Det\cap L$ is equal to the union
\[
Det\cap L=\overline{p((E\backslash q^{-1}(0))\cap L\times V^{\lor})}=\overline{p((E\cap L\times V^{\lor})\backslash Q_{[n]})}=\underset{F\in\mathcal{L}\backslash[n]}{\cup}\overline{p(Q_{F})},
\]
where $\mathcal{L}\backslash[n]$ is the lattice of flats without
the maximal element $[n]$. 
\end{proof}
Recall that a tuple of subspaces $(L_{1},...,L_{n})$ is independent
if the defects of all subtuples are non-negative. According to the
Minkowski theorem, this is equivalent to the existence of $n$ linearly
independent vectors $\upsilon_{i}\in L_{i}$.
\begin{cor}
For a dependent tuple $(L_{1},...,L_{n})$ of vector subspaces, the
subspace $L$ lies in the determinant hypersurface $Det$.
\end{cor}

\begin{thm}
\label{Theorem. Intersection of the determinant with a subspace}For
an irreducible BK-tuple $(L_{1},...,L_{n})$ of vector subspaces,
the intersection of the determinant hypersurface $Det$ with the subspace
$L$ is a variety.
\end{thm}

\begin{proof}
The equation $det\,\varphi|_{L}$ is not trivial, because irreducible
tuples of subspaces are independent. Then every matrix from the intersection
$Det\cap L$ has at least one eigenvector corresponding to eigenvalue
zero, and the dimension of the projection is bounded by $dim\,p(Q_{F})\leq dim\,Q_{F}-1=dim\,L-\delta(F)-1.$
Since the tuple $(L_{1},...,L_{n})$ is irreducible, the defects of
proper flats are positive, and $dim\,p(Q_{F})<dim\,L-1$. Only for
the empty BK-subtuple $(\varnothing)$, the inequality $dim\,p(Q_{\varnothing})\leq dim\,L-1$
is not strict, and there are no other candidates between strata to
have codimension one in $L$. Since the intersection $Det\cap L$
is nonempty and has codimension one, the equality holds $codim_{L}\,p(Q_{\varnothing})=1$.
Since the stratum $Q_{\varnothing}$ is a variety, the projection
$p(Q_{\varnothing})$ is a variety. The closure of $p(Q_{\varnothing})$
contains projections of all other strata $Q_{F}$ and equals the intersection
$Det\cap L$.
\end{proof}
\begin{rem}
For a BK-tuple $(L_{1},...,L_{n})$ of vector subspaces, it is possible
to show that the number of components for the intersection $Det\cap L$
equals the number of elements in the Birkhoff poset for BK-subtuples
\cite{pokidkin_combinatorics_2025}. This completes the description
of the intersection of the determinant hypersurface with a row-generated
subspace.
\end{rem}


\section{\label{Section. Esterov irreducibility theorem}Esterov irreducibility
theorem}

Denote by $\mathbb{K}_{\mathscr{A}}$ the space of polynomial systems
over an algebraically closed field $\mathbb{K}$. We search for solutions
of polynomial systems $\Phi_{\mathscr{A}}$ in the split torus $T(\mathscr{A})=\langle\mathscr{A}\rangle^{\vee}\otimes\mathbb{K}^{\times}$,
$\langle\mathscr{A}\rangle=M$. In the affine space $\mathbb{K}_{\mathscr{A}}$,
the $\mathscr{A}$\textit{-discriminant} is the Zariski closure of
the set of polynomial systems with a degenerate root in the split
torus $T(\mathscr{A})$. Denote by $\mathfrak{c}(\mathscr{A})$ the
number of sets in the tuple $\mathscr{A}$.
\begin{thm}
\label{Theorem. Esterov's Conjecture}For an irreducible BK-tuple
$\mathscr{A}$, the discriminant $D_{\mathscr{A}}$ is irreducible
in $\mathbb{K}_{\mathscr{A}}$.
\end{thm}

\begin{proof}
A tuple $\mathscr{A}$ defines a vector subspace of matrices $Mat_{\mathscr{A}}=\underset{A\in\mathscr{A}}{\oplus}(\langle A\rangle\otimes\mathbb{K})$
in the space of matrices $Mat_{dim\langle\mathscr{A}\rangle,\mathfrak{c}(\mathscr{A})}$.
For the irreducible BK-tuple $\mathscr{A}$, the intersection $Y$
of the determinant hypersurface $Det$ with the subspace $Mat_{\mathscr{A}}$
is a variety by Theorem \ref{Theorem. Intersection of the determinant with a subspace}.

With a finite set $A$ and a tuple $\mathscr{A}$, we associate matrices
$\tau_{A}=(a,\;a\in A)$ and $\tau_{\mathscr{A}}=\underset{A\in\mathscr{A}}{\oplus}\tau_{A}$,
with integer coefficients and ranks $rk\,\tau_{A}=dim\,\langle A\rangle,$
$rk\,\tau_{\mathscr{A}}=dim\,Mat_{\mathscr{A}}$. The matrix $\tau_{\mathscr{A}}$
defines the linear map from $\mathbb{K}_{\mathscr{A}}$ to $Mat_{\mathscr{A}}$
by the formula $\tau_{\mathscr{A}}(\lambda_{\mathscr{A}})=(\tau_{A}(\lambda_{A}),\;A\in\mathscr{A})$,
where $\tau_{A}(\lambda_{A})=\underset{a\in A}{\sum}\lambda_{a}\,a$.
This means that the preimage $K=\tau_{\mathscr{A}}^{-1}(Y)$ is a
variety in $\mathbb{K}_{\mathscr{A}}$ as a trivial vector bundle
of rank $dim\,\mathbb{K}_{\mathscr{A}}-dim\,Mat_{\mathscr{A}}$.

We can construct the following sequence of maps:
\[
\begin{array}{ccccccc}
T(\mathscr{A})\times\mathbb{K}_{\mathscr{A}} & \overset{\sigma_{\mathscr{A}}}{\twoheadrightarrow} & \mathbb{K}_{\mathscr{A}} & \overset{\tau_{\mathscr{A}}}{\rightarrow} & Mat_{\mathscr{A}} & \overset{det}{\rightarrow} & \mathbb{K},\end{array}
\]
where $\sigma_{\mathscr{A}}(x,c)=(c_{a}x^{a})_{a\in A\in\mathscr{A}}$
is the torus action. Notice that the preimage $W=\sigma_{\mathscr{A}}^{-1}(K)$
is a variety that is isomorphic to $T(\mathscr{A})\times K$. Indeed,
the isomorphism is defined by the regular maps $W\stackrel[\iota]{\sigma}{\rightleftarrows}T(\mathscr{A})\times K$
such that $\iota(x,c_{a})=(x,\frac{c_{a}}{x^{a}})$ and $\sigma(x,c_{a})=(x,c_{a}x^{a})$
for every element $a\in A\in\mathscr{A}.$

The sequence of maps defines the equation for a root $x$ of a polynomial
system $c=(c_{a})_{a\in A\in\mathscr{A}}$ to be a singular point:
$det\circ\tau_{\mathscr{A}}\circ\sigma_{\mathscr{A}}(c,x)=det\,||\underset{a\in A}{\sum}c_{a}x^{a}a||_{A\in\mathscr{A}}=0.$
This polynomial is irreducible because $W$ is a variety.

Consider the vector subspaces $\Pi_{A}$ in $\mathbb{K}_{\mathscr{A}}$
defined by the linear equations $\underset{a\in A}{\sum}c_{a}=0$,
and their intersection $\Pi=\underset{A\in\mathscr{A}}{\cap}\Pi_{A}$.
Denote by $c_{A}$ the coefficient corresponding to zero in a finite
set $A$. In each linear equation, we can express the variable $c_{A}=-\underset{a\in A\backslash\{0\}}{\sum}c_{a}.$
Notice that the determinant $det\,||\underset{a\in A}{\sum}c_{a}a||_{A\in\mathscr{A}}=0$
does not depend on the coefficients $c_{A},$ $A\in\mathscr{A}$.
Then the intersection of the variety $K$ with the subspace $\Pi$
is a variety $X=K\cap\Pi$ defined by the same determinant. To see
this, we ensure that the coordinate algebra $\mathbb{K}[X]$ is an
integral domain:
\[
\mathbb{K}[X]=\frac{\mathbb{K}[c_{a},\,a\in A\in\mathscr{A}]}{\langle\underset{a\in A}{\sum}c_{a},\;A\in\mathscr{A};\;det\,||\underset{a\in A}{\sum}c_{a}a||_{A\in\mathscr{A}}\rangle}=\frac{\mathbb{K}[c_{a},\,a\in A\backslash\{0\},\,A\in\mathscr{A}]}{\langle det\,||\underset{a\in A}{\sum}c_{a}a||_{A\in\mathscr{A}}\rangle}.
\]

Thus the preimage $Z=\sigma_{\mathscr{A}}^{-1}(X)$ is a variety,
isomorphic to the direct product $T(\mathscr{A})\times X$ by the
same isomorphisms $\iota$ and $\sigma$. Notice that each preimage
$\sigma_{\mathscr{A}}^{-1}(\Pi_{A})$ equals the zero locus of the
equation $f_{A}(x)=0$ in $T(\mathscr{A})\times\mathbb{K}_{\mathscr{A}}$.
Hence the variety $Z$ projects to the discriminant $D_{\mathscr{A}}$
by the map $T(\mathscr{A})\times\mathbb{K}_{\mathscr{A}}\overset{\pi}{\rightarrow}\mathbb{K}_{\mathscr{A}}$,
$\pi(x,c)=(c)$:
\begin{align*}
Z=\sigma_{\mathscr{A}}^{-1}(X) & =\sigma_{\mathscr{A}}^{-1}(K\cap\Pi)=\sigma_{\mathscr{A}}^{-1}(K)\cap\sigma_{\mathscr{A}}^{-1}(\Pi)=W\underset{A\in\mathscr{A}}{\cap}\sigma_{\mathscr{A}}^{-1}(\Pi_{A})=\\
 & =V(det\,||\underset{a\in A}{\sum}c_{a}x^{a}a||_{A\in\mathscr{A}})\cap\underset{A\in\mathscr{A}}{\cap}V(f_{A}(x)).
\end{align*}
Therefore, the discriminant $D_{\mathscr{A}}$ is a variety as a projection
of a variety.
\end{proof}
\begin{defn}
A linear/nonlinear irreducible BK-tuple is called a \textit{lir}/\textit{nir}.
\end{defn}

\begin{rem}
Linear BK-tuples are the only tuples with unit mixed volume \cite{cattani_mixed_2013,esterov_systems_2015}.
Discriminants for linear tuples were considered in the works \cite{esterov_galois_2019,borger_defectivity_2020}.
\end{rem}

\begin{thm}
\label{Theorem. Codimension of discriminants}For a lir/nir $\mathscr{A}$,
the discriminant $D_{\mathscr{A}}$ is a variety of codimension two/one
in the space of polynomial systems\textup{ $\mathbb{C}_{\mathscr{A}}$.}
\end{thm}

\begin{proof}
If the tuple $\mathscr{A}$ is a nir, then the discriminant $D_{\mathscr{A}}$
has a codimension one component in $\mathbb{C}_{\mathscr{A}}$ according
to \cite{esterov_galois_2019}. By Theorem \ref{Theorem. Esterov's Conjecture},
this component of codimension one is unique, and the discriminant
$D_{\mathscr{A}}$ is a variety.

If the tuple $\mathscr{A}$ is a lir, then the space $\mathbb{C}_{\mathscr{A}}$
corresponds to a space of coefficients for systems of linear equations.
In that case, the discriminant is the intersection of two determinant
hypersurfaces. By Theorem \ref{Theorem. Esterov's Conjecture}, this
intersection is irreducible. We can use the polynomial equations for
the determinants to parameterize the intersection and conclude that
the discriminant has codimension two.
\end{proof}
\begin{cor}
For a lir/nir $\mathscr{A}$, the singular locus of a generic polynomial
system from the discriminant $D_{\mathscr{A}}$ has dimension one/zero.
\end{cor}

\begin{proof}
Consider the dominant map $Z\overset{\pi}{\rightarrow}D_{\mathscr{A}}$
from the proof of Theorem \ref{Theorem. Esterov's Conjecture}. For
a system $\Phi\in D_{\mathscr{A}},$ the fiber $\pi^{-1}(\Phi)=\Phi\times Sing(\Phi)$
is isomorphic to the singular locus of the system $\Phi.$ Using Theorem
\ref{Theorem. Codimension of discriminants}, the dimension of the
generic fiber $\pi^{-1}(\Phi)$ equals 
\begin{align*}
dim\,\pi^{-1}(\Phi) & =dim\,Z-dim\,D_{\mathscr{A}}=\begin{cases}
1, & \text{if the tuple }\mathscr{A}\text{ is a lir,}\\
0, & \text{if the tuple }\mathscr{A}\text{ is a nir.}
\end{cases}
\end{align*}
\end{proof}
\begin{cor}
\label{Corollary. Irreducible BK =00003D> three types of discriminant are equal}For
an irreducible BK-tuple $\mathscr{A}$, the $\mathscr{A}$-discriminant,
the Cayley discriminant, and the mixed discriminant form the same
hypersurface in $\mathbb{C}_{\mathscr{A}}$ \textup{(see \cite{cattani_mixed_2013,esterov_galois_2019})}.
\end{cor}


\section{\label{Section. Combinatorial Review}Combinatorics behind discriminants}

This section provides an overview of combinatorial results concerning
sublattice configurations, tuples of finite sets, and the mixed volume.
For a sublattice tuple $\mathsf{n}=(S_{1},...,S_{k})$ from a lattice
$M$, the notion of irreducibility is motivated by questions from
algebraic geometry. For the simple case of a tuple with $n$ sublattices
in $\mathbb{Z}^{n},$ there is a partition of a reducible tuple into
subtuples corresponding to irreducible ones. The partition of the
tuple $\mathsf{n}$ is defined by the realizable polymatroid on $\mathsf{n}$.
This combinatorial decomposition is a shadow of the decomposition
of an $\mathscr{A}$-discriminant for polynomial systems into irreducible
components.
\begin{defn}
The \textit{saturation} of a sublattice $S$ is the maximal sublattice
$\overline{S}$ containing $S$ with the same dimension as $S$.

The \textit{linear span} $\langle\mathsf{n}\rangle$ of a sublattice
tuple $\mathsf{n}$ is the minimal sublattice containing all sublattices
from the tuple. The \textit{cardinality} $\mathfrak{c}(\mathsf{n})$
is the number of sets in the tuple $\mathsf{n}$, and the \textit{defect}
of $\mathsf{n}$ is the difference $\delta(\mathsf{n})=dim\,\langle\mathsf{n}\rangle-\mathfrak{c}(\mathsf{n}).$

A tuple is \textit{essential} if every proper subtuple has a strictly
greater defect than the whole tuple. A tuple is\textit{ independent}
if the defects of all subtuples are non-negative. An independent tuple
is \textit{irreducible} if the defects of all proper subtuples are
positive. A \textit{BK-tuple} is an independent tuple with zero defect.

For tuples $\mathsf{n}$ and $\mathscr{\mathsf{k}}$ in a lattice
$M,$ the \textit{quotient tuple} is the projection of the complement
subtuple $\mathsf{n}/\mathsf{k}=\pi(\mathsf{n}\backslash\mathsf{k}),$
where $\pi:M\rightarrow M/\overline{\langle\mathsf{k}\rangle}$.
\end{defn}

The notion of a BK-tuple arose from the Kouchnirenko-Bernstein theorem
\cite{bernshtein_number_1975}.
\begin{prop}
\textup{\cite{pokidkin_combinatorics_2025}} For a reducible BK-tuple
$\mathsf{n}$ and a BK-subtuple $\mathsf{k}$,

1) linear spans of irreducible BK-subtuples do not intersect one another
except for the origin;

2) the quotient tuple $\mathsf{n}/\mathsf{k}$ is a BK-tuple;

3) BK-subtuples of the quotient $\mathsf{n}/\mathsf{k}$ are in bijection
with BK-subtuples of $\mathsf{n}$ containing $\mathsf{k}$;

4) quotient tuples satisfy the isomorphism $\mathsf{n}/\mathsf{k}\cong\frac{\mathsf{n}/\mathsf{l}}{\mathsf{k}/\mathsf{l}}$,
where $\mathsf{l}$ is a BK-subtuple of $\mathsf{k}$.
\end{prop}

\begin{defn}
An \textit{order ideal} $\mathtt{I}$ in a poset $\mathtt{P}$ is
a subposet such that if $\beta\in\mathtt{I}$ and $\alpha\leq\beta$,
then $\alpha\in\mathtt{I}$. For an element $\alpha$ in a poset,
the \textit{principal order ideal }$(\alpha)$ is the order ideal
of all elements that are not greater than $\alpha.$

The dual notions are \textit{order filter} and \textit{principal order
filter} $[\alpha]$ (inverse all $\leq$).
\end{defn}

Since intersections and unions of BK-subtuples are BK-subtuples in
a BK-tuple $\mathsf{n}$ \cite{pokidkin_combinatorics_2025}, the
set of BK-subtuples forms a distributive lattice $L$ by inclusion.
By the fundamental theorem for distributive lattices \cite{stanley_enumerative_2011},
there exists a poset $\mathtt{P}$ such that its lattice of order
ideals is isomorphic to $L$.
\begin{defn}
A \textit{filtration} of a tuple $\mathsf{n}$ is an increasing family
of subtuples $F_{0}\mathsf{n}\hookrightarrow F_{1}\mathsf{n}\hookrightarrow...\hookrightarrow F_{m}\mathsf{n}=\mathsf{n}.$
A filtration is a \textit{BK-filtration} if all quotients $F_{j}\mathsf{n}/F_{j-1}\mathsf{n}$
are BK-tuples. A BK-filtration is \textit{maximal} if all quotients
$F_{j}\mathsf{n}/F_{j-1}\mathsf{n}$ are irreducible.
\end{defn}

\begin{prop}
\label{Corollary. Filtration of BK-tuple with MV>0}For a reducible
BK-tuple $\mathsf{n}$, the number of subtuples in a maximal BK-filtration
equals the number of elements in the Birkhoff poset. Moreover, there
exist linear isomorphisms between successive quotients $F_{j}\mathsf{n}/F_{j-1}\mathsf{n}$
and $F_{j}'\mathsf{n}/F_{j-1}'\mathsf{n}$, for any given maximal
BK-filtrations $F_{\bullet}\mathsf{n}$ and $F'_{\bullet}\mathsf{n}.$
\end{prop}

This Birkhoff poset $\mathtt{P}$ defines a partition of the BK-tuple
$\mathsf{n}$: every element $\alpha$ corresponds to some subtuple
$\mathsf{k}_{\alpha}$ of $\mathsf{n}$, and every order ideal $\mathtt{I}$
of $\mathtt{P}$ corresponds to some BK-subtuple $\mathtt{k_{I}}=\underset{\alpha\in\mathtt{I}}{\sqcup}\mathsf{k}_{\alpha}.$ 
\begin{thm}
\textup{\label{Theorem. Poset partition of a reducible BK-tuple }
\cite{pokidkin_combinatorics_2025}} A reducible BK-tuple $\mathsf{n}$
admits the unique partition $\mathsf{n}=\underset{\alpha\in\mathtt{P}}{\sqcup}\mathsf{k}_{\alpha}$
such that the subtuples $\hat{\mathsf{k}}_{\alpha}=\mathsf{k}_{(\alpha)}/\mathsf{k}_{(\alpha)\backslash\alpha}$
are irreducible BK-tuples for every element $\alpha$ of the poset
$\mathtt{P}$.
\end{thm}

We use these results to construct partitions of tuples of finite sets.
The \textit{affine linear span} of a finite set $A$ in a lattice
$M$ is the affine sublattice $\langle A\rangle$ generated by this
set. Using a shift, we can always ensure that the finite set $A$
contains the origin, and its linear span is a sublattice in $M$.
Then a tuple of finite sets $\mathscr{A}=(A_{1},...,A_{m})$ with
the common origin generates a sublattice tuple that can be considered
as a realizable polymatroid.

General polymatroids admit similar characterizations and terminology
as sublattice tuples. For the sake of simplicity and because there
is no necessity, we do not provide results about polymatroids. Instead,
we present results for tuples of sublattices from a more general theory.
However, it was the study of discriminants that prompted results in
the polymatroid theory and inspired the work \cite{pokidkin_combinatorics_2025}.

We characterize a tuple of finite sets via the generated sublattice
tuple.

By the mixed volume $\mathrm{MV}_{M}(\mathscr{A})$ of a tuple of
finite sets $\mathscr{A}$ in a lattice $M$, we mean the mixed volume
of the corresponding convex hulls of finite sets. We highlight that
a BK-tuple $\mathscr{A}$ has a positive mixed volume in the sublattice
$\langle\mathscr{A}\rangle$ by Minkowski's theorem (see Theorem 8
\cite{khovanskii_newton_2016}). Another significant result establishes
the compatibility of the poset partition of a reducible BK-tuple of
finite sets with its mixed volume:
\begin{thm}
\label{Theorem. BK-subtuple and mixed volume}For tuples $\mathscr{B}\subset\mathscr{A}$
with zero defect in a lattice $M$, the mixed volume decomposes 
\[
\mathrm{MV}_{M}(\mathscr{A})=\mathrm{MV}_{\overline{\langle\mathscr{B}\rangle}}(\mathscr{B})\,\mathrm{MV}_{M/\overline{\langle\mathscr{B}\rangle}}(\mathscr{A}/\mathscr{B}).
\]
\end{thm}

\begin{proof}
See Lemma 4 \cite{steffens_mixed_2010} for a geometric proof, Lemma
3.12 \cite{dandrea_poisson_2015} for an analytic proof and Theorem
1.10 \cite{esterov_galois_2019} for an algebraic proof.
\end{proof}
\begin{cor}
\textup{\label{Corollary. Unique decomposition of a BK-tuple}}A reducible
BK-tuple $\mathscr{A}$ admits the unique partition $\mathscr{A}=\underset{\alpha\in\mathtt{P}_{\mathscr{A}}}{\sqcup}\mathscr{B}_{\alpha}$
such that tuples $\hat{\mathscr{B}}_{\alpha}=\mathscr{B}_{(\alpha)}/\mathscr{B}_{(\alpha)\backslash\alpha}$
are irreducible BK-tuples for every element $\alpha$ of the poset
$\mathtt{P}_{\mathscr{A}}.$
\end{cor}

For a BK-tuple, BK-subtuples form a distributive lattice. Hence it
is possible to choose a basis in the ambient lattice such that every
BK-subtuple lies in a coordinate sublattice by Proposition 7.1 \cite{polishchuk_quadratic_2005}.

For a dependent tuple of finite sets, independent subtuples are independent
subsets of the induced matroid from the realizable polymatroid \cite{pokidkin_combinatorics_2025}.
We will need the subtuple that is the maximal cycle of the induced
matroid to characterize the discriminant.

\section{\label{Section. Discriminants-of-polynomial}Discriminants of polynomial
systems}

For a group lattice $N,$ the \textit{dual lattice} is $M=N^{\vee}=\mathrm{Hom}_{\mathbb{Z}}(N,\mathbb{Z}),$
and the \textit{algebraic torus} is $T(N)=N\otimes_{\mathbb{Z}}\mathbb{C}^{\times}.$
The \textit{character group}\textbf{ }is $\mathrm{Hom}_{\mathbb{Z}}(T(N),\mathbb{C}^{\times})$,
and its elements are called \textit{characters} $\chi_{a}$ \cite{fulton_introduction_1993,cox_toric_2011}.
This character group is isomorphic to the dual lattice $M$. For a
tuple of finite sets $\mathscr{A}\subset M$, we have defined the
space of polynomial systems $\mathbb{C}_{\mathscr{A}}$, and every
polynomial system $\Phi$ has solutions in the torus $T(N).$

For a split short exact sequence of lattices $0\rightarrow N'\rightarrow N\overset{\pi}{\rightarrow}N''\rightarrow0$,
we have the corresponding split short exact sequences for the dual
lattices $0\rightarrow M''\overset{i}{\rightarrow}M\rightarrow M'\rightarrow0$,
the character groups and the tori. The projection of lattices $\pi$
induces the \textit{pushforward} $\pi^{*}=\pi\otimes\mathbb{C}^{\times}$
of tori and the \textit{pullback} $\pi_{*}=\mathrm{Hom}\,(\pi^{*},\mathbb{C}^{\times})$
of character groups, $\pi_{*}(\chi_{a})=\chi_{a}\circ\pi^{*}=\chi_{i(a)},$
for a character $\chi_{a}$ from $N''$. 

A monomorphism of dual lattices $M''\overset{i}{\rightarrow}M$ provides
an equality for the discriminants: $\pi_{*}(D_{\mathscr{A}})=D_{i(\mathscr{A})}$.
Hence an isomorphism of dual lattices $M''\overset{i}{\rightarrow}M$
leads to an isomorphism of discriminants: $D_{\mathscr{A}}\cong D_{i(\mathscr{A})}$.
In particular, the discriminants $D_{\mathscr{A}}\cong D_{g\mathscr{A}}$
are isomorphic for every element $g$ of the affine general linear
group $\mathsf{AGL}(n,\mathbb{Z})$. Also, if there exists a one-to-one
correspondence between sets from tuples $\mathscr{A}$ and $\mathscr{B}$
such that each set $A\in\mathscr{A}$ is a translation of a set $B\in\mathscr{B}$,
then the discriminants are equal, $D_{\mathscr{A}}=D_{\mathscr{B}}$. 

These observations allow us to choose convenient tuples for subsequent
proofs and to use the results of Section \ref{Section. Combinatorial Review}.
In the sequel, we assume that the linear span of the tuple $\mathscr{A}$
equals the dual lattice, $\langle\mathscr{A}\rangle=M$. For the torus
$T(N)$ we use the notation $T(\mathscr{A})$.

Consider a split short exact sequence of dual lattices, $0\rightarrow M''\rightarrow M\rightarrow M'\rightarrow0$.
Lattices are reflexive $\mathbb{Z}$-modules (there is an isomorphism
$N\cong\mathrm{Hom}\,(\mathrm{Hom}\,(N,\mathbb{Z}),\mathbb{Z})$)
as finitely generated free $\mathbb{Z}$-modules. This means that
the lattice $\mathrm{Hom}\,(M'',\mathbb{Z})$ is isomorphic to some
lattice $N''$ such that $M''=\mathrm{Hom}\,(N'',\mathbb{Z})$. Let
us compute the lattice $N''$. Notice that the dual sublattice $M''$
naturally corresponds to the sublattice $N'=M''^{\perp}=\{n\in N\,|\;m(n)=0\;\forall m\in M''\}\subseteq N.$
From the splitting, the lattice $N''$ can be defined as the quotient
$N/N'=N/M''^{\perp}$.

For a subtuple $\mathscr{B}$ of a tuple $\mathscr{A}$, we search
for solutions of a subsystem from $\mathbb{C}_{\mathscr{B}}$ in the
special torus $T(N/\overline{\langle\mathscr{B}\rangle}^{\bot}),$
denoted as $T(\mathscr{B})$.

A\textbf{ }\textit{cofiltration} $G_{\bullet}N$ is a sequence of
quotients $N=G_{k}N\twoheadrightarrow G_{k-1}N\twoheadrightarrow...\twoheadrightarrow G_{0}N\twoheadrightarrow0.$
There is a bijection between the cofiltrations of a lattice $N$ and
the filtrations of the dual lattice $M$.

Every BK-tuple $\mathscr{A}$ admits a maximal BK-filtration $F_{\bullet}$.
The saturated linear spans of the tuples from the BK-filtration form
a filtration of the dual lattice $F_{\bullet}M,$ $F_{i}M=\overline{\langle F_{i}\mathscr{A}\rangle.}$
Then there exists a cofiltration $G_{\bullet}N$, and, hence, the
cofiltration for the torus $G_{\bullet}T(N)=T(G_{\bullet}N).$ Therefore,
every polynomial system $\Phi\in\mathbb{C}_{\mathscr{A}}$ admits
a BK-filtration $F_{\bullet}\Phi$, and the solutions of that system
admit a cofiltration $G_{\bullet}V(\Phi).$ If the system $\Phi$
is generic, then the projections of finite sets $G_{i-1}V(\Phi)\twoheadrightarrow G_{i}V(\Phi)$
are finite covers with degrees $\mathrm{MV}_{\text{\ensuremath{G_{i-1}N/G_{i}N}}}(F_{i}\mathscr{A}/F_{i-1}\mathscr{A})$
by Theorem \ref{Theorem. BK-subtuple and mixed volume} and by the
Kouchnirenko-Bernstein theorem. For every subsystem $F_{i}(\Phi_{\mathscr{A}})=\Phi_{\mathscr{B}}$,
we will search for solutions in the torus $T(\mathscr{B})=T(G_{i}N).$

\section{\label{Section. BK-multiplicaiton}BK-multiplication}

To compute discriminants for BK-tuples, we build a specific multiplication
between varieties. 

Consider a split short exact sequence of dual lattices $0\rightarrow M''\rightarrow M\overset{\tau}{\rightarrow}M'\rightarrow0$,
a finite set $A\subset M,$ its projection $B=\tau(A),$ and the corresponding
splitting of the algebraic torus $T(N)=T(N'')\times T(N')$. We can
represent a point in the torus $T(N)$ as a pair $(x,y)\in T(N'')\times T(N').$
Then the substitution $f(x,y)\overset{ev_{x_{0}}}{\longrightarrow}f(x_{0},y)$
of a fixed point $x_{0}$ from the torus $T(N'')$ into polynomials
from $\mathbb{C}_{A}$ is a linear projection on $\mathbb{C}_{B}$,
$\mathbb{C}_{A}\overset{ev_{x_{0}}}{\longrightarrow}\mathbb{C}_{B}$.
For a tuple of finite sets $\mathscr{A}\subset M,$ $\mathscr{B}=\tau(\mathscr{A})$,
the substitution $\Phi(x,y)\overset{ev_{x_{0}}}{\longrightarrow}\Phi(x_{0},y)$
of the point $x_{0}$ corresponds to a linear projection $\mathbb{C}_{\mathscr{A}}\overset{ev_{x_{0}}}{\longrightarrow}\mathbb{C}_{\mathscr{B}}$.
\begin{lem}
\label{Lemma. Evaluation bundle is a trivial vector bundle}For any
quasi-affine algebraic set $Y\subset\mathbb{C}_{\mathscr{B}}$, a
point $x$ from the torus $T(N'')$, and the preimage $E=ev_{x}^{-1}(Y)\subseteq\mathbb{C}_{\mathscr{A}}$,
the triple $(E,ev_{x},Y)$ is a trivial vector bundle over $Y$ of
rank $\underset{A\in\mathscr{A}}{\sum}|A|-|\tau(A)|$.
\end{lem}

\begin{defn}
We call the total space $E$ the \textit{evaluation bundle} over $Y$
for the evaluation by a point $x$ from the torus $T(N'')$ and denote
it by $E_{\mathscr{A}}^{x}(Y)=E$.
\end{defn}

\begin{rem}
\label{Remark. Evaluation bundle}For tuples of finite sets $\mathscr{B}\subset\mathscr{A}$
and a point $x$ from the torus $T(\mathscr{B})$, the evaluation
bundle $E_{\mathscr{A}}^{x}(Y)$ is a variety in $\mathbb{C}_{\mathscr{A}\backslash\mathscr{B}}$
if and only if $Y$ is a variety in $\mathbb{C}_{\mathscr{A}/\mathscr{B}}$. 
\end{rem}

\begin{lem}
\label{Lemma. Fiber change}For a chain of BK-tuples $\mathscr{C}\subset\mathscr{B}\subset\mathscr{A}$,
a quasi-affine algebraic set $Y\subset\mathbb{C}_{\mathscr{B}/\mathscr{C}}$
and a point $x\in T(\mathscr{C})$, the following holds:
\[
E_{\mathscr{A}\backslash\mathscr{C}}^{x}\left(Y\times\mathbb{C}_{(\mathscr{A}\backslash\mathscr{B})/\mathscr{C}}\right)=E_{\mathscr{B}\backslash\mathscr{C}}^{x}(Y)\times\mathbb{C}_{\mathscr{A}\backslash\mathscr{B}}.
\]
\end{lem}

\begin{proof}
Notice that the evaluation map $ev_{x}:\mathbb{C}_{\mathscr{A}\backslash\mathscr{C}}\rightarrow\mathbb{C}_{\mathscr{A}/\mathscr{C}}$
splits into two 

$u_{x}:\mathbb{C}_{\mathscr{B}\backslash\mathscr{C}}\rightarrow\mathbb{C}_{\mathscr{B}/\mathscr{C}}$
and $v_{x}:\mathbb{C}_{\mathscr{A}\backslash\mathscr{B}}\rightarrow\mathbb{C}_{(\mathscr{A}\backslash\mathscr{B})/\mathscr{C}}$,
$ev_{x}=u_{x}\oplus v_{x}.$ 

Then we have $ev_{x}^{-1}(Y\times\mathbb{C}_{(\mathscr{A}\backslash\mathscr{B})/\mathscr{C}})=u_{x}^{-1}(Y)\times v_{x}^{-1}(\mathbb{C}_{(\mathscr{A}\backslash\mathscr{B})/\mathscr{C}}).$
\end{proof}
\begin{thm}
\textup{\label{Theorem. Reformulation of Kouchnirenko-Bernstein theorem}(Kouchnirenko-Bernstein,
\cite{bernshtein_number_1975})} For a tuple $\mathscr{A}$ of $n$
finite sets in an $n$-dimensional lattice $M$, there exists an open
subset $U$ in $\mathbb{C}_{\mathscr{A}}$ such that the set of solutions
for every polynomial system $\Phi\in U$ consists of exactly $\mathrm{MV}_{M}(\mathscr{A})$-points.
\end{thm}

The complement $\mathbb{C}_{\mathscr{A}}\backslash U$ is a \textit{bifurcation
divisor} according to Esterov's work \cite{esterov_discriminant_2013}.
Discriminants of different types lie in the bifurcation divisor.
\begin{defn}
\label{Definition. BK-multiplication}For BK-tuples of finite sets
$\mathscr{B}\subset\mathscr{A},$ the \textit{BK-multiplication} $X\circ Y\subseteq\mathbb{C}_{\mathscr{A}}$
of algebraic sets $X\subset\mathbb{C}_{\mathscr{B}}$ and $Y\subset\mathbb{C}_{\mathscr{A}/\mathscr{B}}$
is called the quasi-affine set $\{\Phi\times\underset{x\in V(\Phi)}{\cup}ev_{x}^{-1}(Y)\,|\;\Phi\in X\}\subset\mathbb{C}_{\mathscr{A}}$,
where $V(\Phi)$ is a set of zeroes for a polynomial system $\Phi\in X$
in the torus $T(\mathscr{B})$ and the preimage $ev_{x}^{-1}(Y)=E_{\mathscr{A}\backslash\mathscr{B}}^{x}(Y)$
is an evaluation bundle in $\mathbb{C}_{\mathscr{A}\backslash\mathscr{B}}$
over $Y$ for the evaluation by a root $x\in T(\mathscr{B})$ of the
polynomial system $\Phi$. Denote by $X\bullet Y$ the Zariski closure
of the BK-multiplication $X\circ Y.$
\end{defn}

\begin{rem}
1) By the Kouchnirenko-Bernstein Theorem \ref{Theorem. Reformulation of Kouchnirenko-Bernstein theorem},
the generic fiber in the BK-multiplication is a union of $\mathrm{MV}_{\overline{\langle\mathscr{B}\rangle}}(\mathscr{B})$
different trivial evaluation bundles over $Y$.

2) Denote by $Z_{\mathscr{B}}$ the set of polynomial systems in $\mathbb{C}_{\mathscr{B}}$
with an empty set of solutions. Then the contribution $\Phi\times\underset{x\in V(\Phi)}{\cup}ev_{x}^{-1}(Y)$
is empty for every polynomial system $\Phi\in Z_{\mathscr{B}}$. That
is why we take the algebraic closure of the multiplication.
\end{rem}

\begin{lem}
\label{Lemma. BK-multiplication for C}The following holds: $\mathbb{C}_{\mathscr{A}}=\mathbb{C}_{\mathscr{B}}\bullet\mathbb{C}_{\mathscr{A}/\mathscr{B}}$.
\end{lem}

\begin{proof}
Every system $\Phi$ from $\mathbb{C}_{\mathscr{B}}\backslash Z_{\mathscr{B}}$
has at least one solution $x$. Hence every fiber equals the same
vector space $\underset{x\in V(\Phi)}{\cup}ev_{x}^{-1}(\mathbb{C}_{\mathscr{A}/\mathscr{B}})=\mathbb{C}_{\mathscr{A}\backslash\mathscr{B}}$.
Therefore, the algebraic closure of the set $\{\Phi\times\mathbb{C}_{\mathscr{A}\backslash\mathscr{B}}\,|\;\Phi\in\mathbb{C}_{\mathscr{B}}\backslash Z_{\mathscr{B}}\}$
coincides with the space $\mathbb{C}_{\mathscr{A}}.$ This lemma is
a shadow of the mixed volume decomposition in Theorem \ref{Theorem. BK-subtuple and mixed volume}.
\end{proof}
\begin{cor}
\label{Corollay. BK-multiplicaiton for XoC}Every variety $X$ from
$\mathbb{C}_{\mathscr{B}}\backslash Z_{\mathscr{B}}$ satisfies the
equality: $X\bullet\mathbb{C}_{\mathscr{A}/\mathscr{B}}=\overline{X}\times\mathbb{C}_{\mathscr{A}\backslash\mathscr{B}}$.
\end{cor}

\begin{lem}
\label{Lemma. The set of solutions is a variety}For a tuple $\mathscr{A},$
the set of solutions $V(\Phi_{\mathscr{A}}(x))$ is a variety in $T(\mathscr{A})\times\mathbb{C}_{\mathscr{A}}.$
\end{lem}

\begin{proof}
Let $\Pi=\underset{A\in\mathscr{A}}{\cap}\Pi_{A}$ be an intersection
of hyperplanes $\Pi_{A}$ in $\mathbb{C}_{\mathscr{A}}$, defined
by the equations $\underset{a\in A}{\sum}c_{a}=0$ for each set $A\in\mathscr{A}$.
Consider the map $T(\mathscr{A})\times\mathbb{C}_{\mathscr{A}}\overset{\sigma_{\mathscr{A}}}{\twoheadrightarrow}\mathbb{C}_{\mathscr{A}}$
such that $\sigma_{\mathscr{A}}(x,c)=(c_{a}x^{a})_{a\in A\in\mathscr{A}}.$
Notice that the preimage $W=\sigma_{\mathscr{A}}^{-1}(\Pi)$ is a
variety isomorphic to the product $T(\mathscr{A})\times\Pi.$ Indeed,
the isomorphism is defined by the regular maps $W\stackrel[\iota]{\sigma}{\rightleftarrows}T(\mathscr{A})\times\Pi$
such that $\iota(c_{a})=(\frac{c_{a}}{x^{a}})$ and $\sigma(c_{a})=(c_{a}x^{a})$
for every $a\in A\in\mathscr{A}$ and $x\in T(\mathscr{A})$ ($\iota$
and $\sigma$ are identity on other coordinates). Moreover, the preimage
of each hyperplane $\sigma_{\mathscr{A}}^{-1}(\Pi_{A})$ equals the
zero locus of the equation $f_{A}(x)=0$ in $T(\mathscr{A})\times\mathbb{C}_{\mathscr{A}}$.
Hence the variety $W$ equals the set of solutions $V(\Phi_{\mathscr{A}}(x))$.
\end{proof}
\begin{thm}
\label{Theorem. BK-multiplication of varieties}For a variety $Y\subset\mathbb{C}_{\mathscr{A}/\mathscr{B}},$
the BK-multiplication $\mathbb{C}_{\mathscr{B}}\circ Y$ is a variety.
\end{thm}

\begin{proof}
By Lemma \ref{Lemma. The set of solutions is a variety}, the set
of solutions $V(\Phi_{\mathscr{B}}(x))$ is a variety $W$ in $T(\mathscr{B})\times\mathbb{C}_{\mathscr{B}}$
for a BK-subtuple $\mathscr{B}$. Consider the evaluation map 
\[
T(\mathscr{B})\times\mathbb{C}_{\mathscr{A}}\overset{ev}{\twoheadrightarrow}T(\mathscr{B})\times\mathbb{C}_{\mathscr{B}}\times\mathbb{C}_{\mathscr{A}/\mathscr{B}}
\]
such that $ev(x,\Phi_{\mathscr{B}},\Phi_{\mathscr{A}\backslash\mathscr{B}})=(x,\Phi_{\mathscr{B}},ev_{x}(\Phi_{\mathscr{A}\backslash\mathscr{B}}))$.
Notice that the product $W\times Y$ is a variety, and the preimage
$ev^{-1}(W\times Y)$ is a trivial vector bundle over $W\times Y$.
Indeed, for a fixed $x\in T(\mathscr{B}),$ the preimage is defined
by linear equations $\underset{a\in A_{b}}{\sum}c_{a}x^{a}=c_{b}$
for every $b\in B\in\mathscr{A}/\mathscr{B}$ (every set $A$ in $\mathscr{A}\backslash\mathscr{B}$
admits the partition $A=\underset{b\in B}{\sqcup}A_{b}$ for the corresponding
$B\in\mathscr{A}/\mathscr{B}$). Then the projection of the variety
$ev^{-1}(W\times Y)$ on the space $\mathbb{C}_{\mathscr{A}}$ is
a variety that equals the BK-multiplication $\mathbb{C}_{\mathscr{B}}\circ Y$.
\end{proof}
\begin{rem}
The closed BK-multiplication is associative: $(X\bullet Y)\bullet Z=X\bullet(Y\bullet Z)$
for $X\subseteq\mathbb{C}_{\mathscr{C}},$ $Y\subseteq\mathbb{C}_{\mathscr{B}/\mathscr{C}},$
$Z\subseteq\mathbb{C}_{\mathscr{A}/\mathscr{B}}$, and a chain of
BK-tuples $\mathscr{C}\subset\mathscr{B}\subset\mathscr{A}.$

The closed BK-multiplication is distributive in the following sense:
$X\bullet Y\cup X'\bullet Y=(X\cup X')\bullet Y$ and $X\bullet Y\cup X\bullet Y'=X\bullet(Y\cup Y')$
for algebraic sets $X,X'\subseteq\mathbb{C}_{\mathscr{B}}$, $Y,Y'\subseteq\mathbb{C}_{\mathscr{A}/\mathscr{B}}$,
and BK-tuples $\mathscr{B}\subset\mathscr{A}$. The closed BK-multiplication
is commutative $X\bullet Y=Y\bullet X$ only if the complement $\mathscr{A}\backslash\mathscr{B}$
is a BK-tuple.
\end{rem}

\begin{cor}
\label{Corollary. Coherency Relations}\textup{(Coherency Relations)}
For a chain of BK-tuples $\mathscr{C}\subset\mathscr{B}\subset\mathscr{A}$
and quasi-affine algebraic sets $Y_{\mathscr{B}/\mathscr{C}}\subseteq\mathbb{C}_{\mathscr{B}/\mathscr{C}}\backslash Z_{\mathscr{B}/\mathscr{C}}$
and $Y_{\mathscr{A}/\mathscr{B}}\subseteq\mathbb{C}_{\mathscr{A}/\mathscr{B}}$,
the equalities hold:
\begin{align*}
\mathbb{C}_{\mathscr{C}}\bullet(Y_{\mathscr{B}/\mathscr{C}}\times\mathbb{C}_{(\mathscr{A}\backslash\mathscr{B})/\mathscr{C}}) & =\mathbb{C}_{\mathscr{C}}\bullet Y_{\mathscr{B}/\mathscr{C}}\bullet\mathbb{C}_{\mathscr{A}/\mathscr{B}},\\
\mathbb{C}_{\mathscr{B}}\bullet Y_{\mathscr{A}/\mathscr{B}} & =\mathbb{C}_{\mathscr{C}}\bullet\mathbb{C}_{\mathscr{B}/\mathscr{C}}\bullet Y_{\mathscr{A}/\mathscr{B}}.
\end{align*}
\end{cor}

\begin{proof}
By associativity, Lemma \ref{Lemma. BK-multiplication for C} and
Corollary \ref{Corollay. BK-multiplicaiton for XoC}.
\end{proof}

\section{\label{Section. Discriminants for BK-tuples}$\mathscr{A}$-discriminants}
\begin{lem}
\label{Lemma. Linear Algebra}In a vector space $V,$ consider a linearly
dependent set of vectors $\{\upsilon_{i}\}_{i\in I}$ and a linearly
independent subset $\{\upsilon_{j}\}_{j\in J}$ , $J\subset I.$ Then
the set $\{\pi(\upsilon_{k})\}_{k\in I\backslash J}$ is linearly
dependent in $U$ for the projection $\pi:V\rightarrow U=V/\overline{\langle\upsilon_{j}\rangle_{j\in J}}$. 
\end{lem}

\begin{proof}
If we denote by $W$ the kernel of $\pi$, then the vector space $V$
is isomorphic to the direct sum $V\cong U\oplus W.$ Since the set
$\{\upsilon_{j}\}_{j\in J}$ is linearly independent, there exist
constants $c_{i}$ such that $\underset{i\in I}{\sum}c_{i}\upsilon_{i}=0$.
From the isomorphism, we can decompose each vector $\upsilon_{i}=u_{i}+w_{i}.$
Notice that the vectors $\upsilon_{j}$ lie in the subspace $W$ for
$j\in J$, and $u_{j}=0$. Therefore, the projected subset is linearly
dependent: $\underset{k\in I\backslash J}{\sum}c_{k}u_{k}=0$, $u_{k}=\pi(\upsilon_{k}).$
\end{proof}
\begin{thm}
\label{Theorem. Discriminatn's splitting}For BK-tuples $\mathscr{B}\subset\mathscr{A}$,
the discriminant $D_{\mathscr{A}}$ equals the union 
\[
D_{\mathscr{A}}=D_{\mathscr{B}}\bullet\mathbb{C}_{\mathscr{A}/\mathscr{B}}\cup\mathbb{C}_{\mathscr{B}}\bullet D_{\mathscr{A}/\mathscr{B}}.
\]
\end{thm}

\begin{proof}
For a generic system $\Phi\in D_{\mathscr{A}}$, there exists a singular
point $(x,y)\in T(\mathscr{A})$ such that vectors $\{df_{A}(x,y)\}_{A\in\mathscr{A}}$
are linearly dependent, and $x\in T(\mathscr{B})$. If the set $\{df_{B}(x)\}_{B\in\mathscr{B}}$
is linearly dependent, then the point $\Phi$ belongs to $D_{\mathscr{B}}\bullet\mathbb{C}_{\mathscr{A}/\mathscr{B}}.$
Otherwise, the set $\{d(ev_{x}f_{C})(y)\}_{C\in\mathscr{A}/\mathscr{B}}$
is linearly dependent by Lemma \ref{Lemma. Linear Algebra}, and the
substitution of the root $x$ into the system $\Phi_{\mathscr{A}\backslash\mathscr{B}}$
gives us the singular system $ev_{x}\Phi_{\mathscr{A}\backslash\mathscr{B}}$
from the discriminant $D_{\mathscr{A}/\mathscr{B}}.$ Then the system
$\Phi$ lies in $\mathbb{C}_{\mathscr{B}}\bullet D_{\mathscr{A}/\mathscr{B}}$.
\end{proof}
The theorem above splits the discriminant $D_{\mathscr{A}}$ into
two parts: $D_{\mathscr{B}}\bullet\mathbb{C}_{\mathscr{A}/\mathscr{B}}$
and $\mathbb{C}_{\mathscr{B}}\bullet D_{\mathscr{A}/\mathscr{B}}$.
Each part can consist of a collection of components. However, the
theorem does not describe inclusions between the parts. The theorem
shows the possibility of having more than one component. Nevertheless,
the result is significant since, with every BK-subtuple $\mathscr{B}$,
we can associate some part $C(\mathscr{B})$ in the discriminant $D_{\mathscr{A}}$
by a sequence of splittings and usage of coherency relations for BK-multiplications.
An irreducible part is called a \textit{stratum}.
\begin{prop}
\label{Proposition. The second strata doesn't lie in the first}The
part $\mathbb{C}_{\mathscr{B}}\bullet D_{\mathscr{A}/\mathscr{B}}$
does not lie in the other $D_{\mathscr{B}}\bullet\mathbb{C}_{\mathscr{A}/\mathscr{B}}$.
\end{prop}

\begin{proof}
The part $\mathbb{C}_{\mathscr{B}}\bullet D_{\mathscr{A}/\mathscr{B}}$
contains the dense subset $(\mathbb{C}_{\mathscr{B}}\backslash B_{\mathscr{B}})\circ D_{\mathscr{A}/\mathscr{B}}$
that does not lie in the closed BK-multiplication $D_{\mathscr{B}}\bullet\mathbb{C}_{\mathscr{A}/\mathscr{B}}.$
\end{proof}
\begin{cor}
\label{Corollary. Component of a poset}Suppose the poset $\mathtt{P}_{\mathscr{A}}$
(see Corollary \ref{Corollary. Unique decomposition of a BK-tuple})
has a few components, and denote by $\mathtt{I}$ one of them; then
the parts $C(\mathscr{B}_{\mathtt{I}})$ and $C(\mathscr{A}\backslash\mathscr{B}_{\mathtt{I}})$
are not contained in each other.
\end{cor}

\begin{proof}
Notice that $\mathtt{I}$ and $\mathtt{P}_{\mathscr{A}}\backslash\mathtt{I}$
are order ideals corresponding to BK-subtuples and use Proposition
\ref{Proposition. The second strata doesn't lie in the first} twice.
\end{proof}
\begin{lem}
\textup{(Separation Lemma)} \label{Lemma. Separation Lemma}For BK-tuples
$\mathscr{C}\subset\mathscr{B}\subset\mathscr{A}$, the corresponding
part $C(\mathscr{B}\backslash\mathscr{C})$ in the discriminant $D_{\mathscr{A}}$
is $\mathbb{C}_{\mathscr{C}}\bullet D_{\mathscr{B}/\mathscr{C}}\bullet\mathbb{C}_{\mathscr{A}/\mathscr{B}}.$
\end{lem}

\begin{proof}
Direct computations with use of Theorem \ref{Theorem. Discriminatn's splitting}
and Corollary \ref{Corollary. Coherency Relations}:
\begin{align*}
D_{\mathscr{A}} & =D_{\mathscr{C}}\bullet\mathbb{C}_{\mathscr{A}/\mathscr{C}}\cup\mathbb{C}_{\mathscr{C}}\bullet(D_{\mathscr{B}/\mathscr{C}}\bullet\mathbb{C}_{\frac{\mathscr{A}/\mathscr{C}}{\mathscr{B}/\mathscr{C}}}\cup\mathbb{C}_{\mathscr{B}/\mathscr{C}}\bullet D_{\frac{\mathscr{A}/\mathscr{C}}{\mathscr{B}/\mathscr{C}}})=\\
 & =D_{\mathscr{C}}\bullet\mathbb{C}_{\mathscr{A}/\mathscr{C}}\cup\mathbb{C}_{\mathscr{C}}\bullet D_{\mathscr{B}/\mathscr{C}}\bullet\mathbb{C}_{\mathscr{A}/\mathscr{B}}\cup\mathbb{C}_{\mathscr{B}}\bullet D_{\mathscr{A}/\mathscr{B}}.
\end{align*}
\end{proof}
\begin{cor}
\label{Corollary. Subtuple from the decomposition corresponds to irreducible stratum}Every
element of the poset $\mathtt{P}_{\mathscr{A}}$ corresponds to some
stratum of $D_{\mathscr{A}}.$
\end{cor}

\begin{proof}
Corollary \ref{Corollary. Unique decomposition of a BK-tuple} provides
the unique decomposition of the tuple $\mathscr{A}$, encoded by a
poset $\mathtt{P}_{\mathscr{A}}$. Every element $\alpha$ of the
poset $\mathtt{P}_{\mathscr{A}}$ corresponds to a subtuple $\mathscr{B}_{\alpha}$
such that the BK-tuple $\hat{\mathscr{B}}_{\alpha}$ is \textit{\emph{irreducible}}.
The corresponding part in the discriminant $D_{\mathscr{A}}$ is $\mathbb{C}_{\mathscr{B}_{(\alpha)\backslash\alpha}}\bullet D_{\hat{\mathscr{B}}_{\alpha}}\bullet\mathbb{C}_{\mathscr{A}/\mathscr{B}_{(\alpha)}}$
by Separation Lemma \ref{Lemma. Separation Lemma}. This part is irreducible
and forms a stratum by Theorems \ref{Theorem. BK-multiplication of varieties}
and \ref{Theorem. Codimension of discriminants} about the irreducibility
of BK-multiplication and the discriminant $D_{\hat{\mathscr{B}}_{\alpha}}$.
\end{proof}
The main goal now is to figure out which strata are not contained
in each other.
\begin{lem}
\label{Lemma. Incomparable elements}For incomparable elements $\alpha,\beta$
of the poset $\mathtt{P}_{\mathscr{A}}$, the corresponding strata
$C(\mathscr{B}_{\alpha})$ and $C(\mathscr{B}_{\beta})$ are not contained
in each other. 
\end{lem}

\begin{proof}
Let us choose the order ideal corresponding to the union of principal
order ideals $(\alpha)\cup(\beta)$ of the poset $\mathtt{P}_{\mathscr{A}}$
and the BK-subtuple $\mathscr{B}=\mathscr{B}_{(\alpha)\cup(\beta)}$.
By Proposition \ref{Proposition. The second strata doesn't lie in the first},
we can explore the strata $C(\mathscr{B}_{\alpha})$ and $C(\mathscr{B}_{\beta})$
in the discriminant $D_{\mathscr{B}}.$ Consider the intersection
of the principal order ideals $(\alpha)\cap(\beta)$ and the BK-subtuple
$\mathscr{C}=\mathscr{B}_{(\alpha)\cap(\beta)}$. By Theorem \ref{Theorem. Discriminatn's splitting},
we can simplify the task by taking the new tuple $\mathscr{B}/\mathscr{C}$
and the discriminant $D_{\mathscr{B}/\mathscr{C}}$. However, the
elements $\alpha,\beta$ are not connected in the poset $(\alpha)\cup(\beta)\backslash((\alpha)\cap(\beta))$.
Therefore, the corresponding strata $C(\mathscr{B}_{\alpha}/\mathscr{C})$
and $C(\mathscr{B}_{\beta}/\mathscr{C})$ do not contain one another
by Corollary \ref{Corollary. Component of a poset}. Consequently,
the strata $C(\mathscr{B}_{\alpha})$ and $C(\mathscr{B}_{\beta})$
are not contained in each other.
\end{proof}
\begin{defn}
In a poset, the \textit{length} of a chain $C$ is the number $\ell(C)$
that is one less than the number of elements in the chain. The \textit{height
}of an element $\alpha$ of a poset is the maximal length of a chain
from the order ideal $(\alpha)$ and is denoted by $h(\alpha)$.

For comparable elements $\alpha\leq\beta$, the \textit{closed interval}
$[\alpha,\beta]$ is an intersection of the principal order ideal
$(\beta)$ with the principal order filter $[\alpha]$.
\end{defn}

\begin{thm}
\label{Theorem. Height Theorem}\textup{(Height Theorem)} If $h(\alpha)<h(\beta)$
for elements $\alpha,\beta$ in the poset $\mathtt{P}_{\mathscr{A}}$,
then the stratum $C(\mathscr{B}_{\beta})$ does not lie in the stratum
$C(\mathscr{B}_{\alpha})$.
\end{thm}

\begin{proof}
For incomparable elements $\alpha,\beta$, use Lemma \ref{Lemma. Incomparable elements}.
For comparable elements $\alpha<\beta$, we reduce the case to the
subposet $[\alpha,\beta]$ - closed interval, and the tuple $\tilde{\mathscr{A}}=\mathscr{B}_{(\beta)}/\mathscr{B}_{(\beta)\backslash[\alpha,\beta]}$
by Separation Lemma \ref{Lemma. Separation Lemma}. Then we apply
calculations from Separation Lemma \ref{Lemma. Separation Lemma}
to the tuple $\mathscr{A}=\tilde{\mathscr{A}}$ and its subtuples
$\mathscr{C}=\tilde{\mathscr{B}}_{\alpha},$ $\mathscr{B}=\mathscr{\tilde{B}}_{[\alpha,\beta]\backslash\beta},$
and use Proposition \ref{Proposition. The second strata doesn't lie in the first}. 
\end{proof}
\begin{lem}
\label{Lemma. Nir subtuples}If the BK-tuple $\hat{\mathscr{B}}_{\alpha}=\mathscr{B}_{(\alpha)}/\mathscr{B}_{(\alpha)\backslash\alpha}$
is a nir, for some element $\alpha$ of the poset $\mathtt{P}_{\mathscr{A}}$,
then the stratum $C(\mathscr{B}_{\alpha})$ is a hypersurface component
in the discriminant $D_{\mathscr{A}}$.
\end{lem}

\begin{proof}
By Height Theorem \ref{Theorem. Height Theorem} and Separation Lemma
\ref{Lemma. Separation Lemma}, we can assume that the element $\alpha$
has zero height. The corresponding stratum is $C(\mathscr{B}_{\alpha})=D_{\mathscr{B}_{\alpha}}\bullet\mathbb{C}_{\mathscr{A}/\mathscr{B}_{\alpha}}$,
and it has codimension one by Theorem \ref{Theorem. Codimension of discriminants}.
The stratum $C(\mathscr{B}_{\alpha})$ has an empty intersection with
$(\mathbb{C}_{\mathscr{B}_{\alpha}}\backslash B_{\mathscr{B}_{\alpha}})\circ D_{\mathscr{A}/\mathscr{B}_{\alpha}}$,
a dense open subset of $C(\mathscr{B}_{\mathtt{P}_{\mathscr{A}}\backslash\alpha})$
of codimension at least one. Therefore, the intersection $C(\mathscr{B}_{\alpha})\cap C(\mathscr{B}_{\mathtt{P}_{\mathscr{A}}\backslash\alpha})$
has codimension at least two, and $C(\mathscr{B}_{\alpha})\cancel{\subset}C(\mathscr{B}_{\mathtt{P}_{\mathscr{A}}\backslash\alpha})$.
\end{proof}
For a reducible BK-tuple $\mathscr{A}$ and an element $\alpha$ of
the poset $\mathtt{P}_{\mathscr{A}},$ the simple BK-subtuple $\mathscr{B}_{(\alpha)}$
is prelinear if and only if the tuple $\hat{\mathscr{B}}_{\alpha}=\mathscr{B}_{(\alpha)}/\mathscr{B}_{(\alpha)\backslash\alpha}$
is linear.
\begin{lem}
\label{Lemma. Linear don't lie in each other}If BK-tuples $\hat{\mathscr{B}}_{\beta}$
are lir for every element $\beta$ from the principal order filter
$[\alpha]$, then the stratum $C(\mathscr{B}_{\alpha})$ is a component
of codimension two in the discriminant $D_{\mathscr{A}}$.
\end{lem}

\begin{proof}
For the principal order filter $[\alpha],$ its complement $\mathtt{P}_{\mathscr{A}}\backslash[\alpha]$
is an order ideal. By Separation Lemma \ref{Lemma. Separation Lemma}
and Theorem \ref{Theorem. Discriminatn's splitting}, we can explore
the discriminant for the new tuple $\tilde{\mathscr{A}}=\mathscr{A}/\mathscr{B}_{\mathtt{P}_{\mathscr{A}}\backslash[\alpha]}$
with the new poset $[\alpha].$ Hence, without loss of generality,
we assume that $\mathtt{P}_{\mathscr{A}}=[\alpha].$ By Height Theorem
\ref{Theorem. Height Theorem}, there are no inclusions $C(\mathscr{B}_{\alpha})\cancel{\supset}C(\mathscr{B}_{\beta})$
for every $\beta>\alpha.$ 

By Separation Lemma \ref{Lemma. Separation Lemma}, the corresponding
stratum is $C(\mathscr{B}_{\alpha})=D_{\mathscr{B}_{\alpha}}\bullet\mathbb{C}_{\mathscr{A}/\mathscr{B}_{\alpha}}$,
and it has codimension two by Theorem \ref{Theorem. Codimension of discriminants}.
The stratum $C(\mathscr{B}_{\alpha})$ has an empty intersection with
$(\mathbb{C}_{\mathscr{B}_{\alpha}}\backslash B_{\mathscr{B}_{\alpha}})\circ D_{\mathscr{A}/\mathscr{B}_{\alpha}}$,
a dense open subset of $C(\mathscr{B}_{[\alpha]\backslash\alpha})$
of codimension two. Therefore, the intersection $C(\mathscr{B}_{\alpha})\cap C(\mathscr{B}_{[\alpha]\backslash\alpha})$
has codimension at least three, and $C(\mathscr{B}_{\alpha})\cancel{\subset}C(\mathscr{B}_{\mathtt{P}_{\mathscr{A}}\backslash\alpha})$.
\end{proof}
\begin{lem}
\label{Lemma. If the upper filer consists of linear tuples then strata codim2 lies in codim1}Consider
BK-tuples $\mathscr{C}\subset\mathscr{B}$ such that $\mathscr{C}$
is linear, and $\mathscr{B}/\mathscr{C}$ is a nir. Then there exists
an inclusion for the strata $C(\mathscr{C})\subset C(\mathscr{B}\backslash\mathscr{C}).$
\end{lem}

\begin{proof}
By Theorem \ref{Theorem. Codimension of discriminants}, the discriminant
$D_{\mathscr{B}/\mathscr{C}}$ is defined by one irreducible polynomial
$P(\Phi_{\mathscr{B}/\mathscr{C}})$ since the quotient tuple $\mathscr{B}/\mathscr{C}$
is a nir. By Remark \ref{Remark. Evaluation bundle}, the evaluation
bundle $E_{\mathscr{B}\backslash\mathscr{C}}^{x}(D_{\mathscr{B}/\mathscr{C}})$
is also irreducible, and it is defined by the polynomial $P_{x}(\Phi_{\mathscr{B}\backslash\mathscr{C}}),$
obtained from the polynomial $P(\Phi_{\mathscr{B}/\mathscr{C}})$
by changing each coefficient $c_{b}$ to the linear polynomial $ev_{x}^{-1}(c_{b})$
for every $b\in C\in\mathscr{B}/\mathscr{C}$ and some fixed point
$x\in T(\mathscr{C})$. Denote by $d$ the degree $deg\,P(\Phi_{\mathscr{B}/\mathscr{C}})$.

Consider a generic point $(\Phi_{\mathscr{C}}^{\circ},\Phi_{\mathscr{B}\backslash\mathscr{C}})\in D_{\mathscr{C}}\times\mathbb{C}_{\mathscr{B}\backslash\mathscr{C}}$.
Since the polynomial system $\Phi_{\mathscr{C}}^{\circ}$ is a degenerate
linear system of more than one variable, the zero locus $V(\Phi_{\mathscr{C}}^{\circ})$
is a vector subspace in $\mathbb{C}(\mathscr{C})=\langle\mathscr{C}\rangle^{\vee}\otimes\mathbb{C}$
of positive dimension. The dimension of $V(\Phi_{\mathscr{C}}^{\circ})$
usually equals one. However, if it is not the case, then choose an
arbitrary linear subspace $U\subseteq V(\Phi_{\mathscr{C}}^{\circ})$.
By definition of the discriminant, in the general case, it is possible
to choose the subspace $U$ intersecting the torus $T(\mathscr{C})$.
If we parameterize $U$ in a natural way, $U=\{at+b,\text{ for some }a,b\in\mathbb{C}(\mathscr{C})\,|\;t\in\mathbb{C}\},$
and expand the brackets in the polynomial $P_{at+b}(\Phi_{\mathscr{B}\backslash\mathscr{C}})$,
then we obtain:
\[
P_{at+b}(\Phi_{\mathscr{B}\backslash\mathscr{C}})=P_{a}(\Phi_{\mathscr{B}\backslash\mathscr{C}})t^{d}+Q_{d-1}(\Phi_{\mathscr{B}\backslash\mathscr{C}})t^{d-1}+...+Q_{1}(\Phi_{\mathscr{B}\backslash\mathscr{C}})t+P_{b}(\Phi_{\mathscr{B}\backslash\mathscr{C}}),
\]
where $Q_{j}(\Phi_{\mathscr{B}\backslash\mathscr{C}})$ are some homogeneous
polynomials of degree $d$. For a fixed point $\Phi_{\mathscr{B}\backslash\mathscr{C}}^{\circ}$,
the polynomial $P_{at+b}(\Phi_{\mathscr{B}\backslash\mathscr{C}}^{\circ})$
of one variable in $t$ has $d$ roots with multiplicities if $P_{a}(\Phi_{\mathscr{B}\backslash\mathscr{C}}^{\circ})\neq0$
and $P_{b}(\Phi_{\mathscr{B}\backslash\mathscr{C}}^{\circ})\neq0.$
Denote by $Y_{x}$ the zero locus $V(P_{x}(\Phi_{\mathscr{B}\backslash\mathscr{C}}))\subset\mathbb{C}_{\mathscr{B}\backslash\mathscr{C}}$
for a fixed point $x\in T(\mathscr{C}).$ Hence there exist $d$ non-zero
solutions with multiplicities of the equation $P_{at+b}(\Phi_{\mathscr{B}\backslash\mathscr{C}})=0$
for every system $\Phi_{\mathscr{B}\backslash\mathscr{C}}\in\mathbb{C}_{\mathscr{B}\backslash\mathscr{C}}\backslash(Y_{a}\cup Y_{b}).$
These roots correspond to the evaluations sending the system $\Phi_{\mathscr{B}\backslash\mathscr{C}}$
to the discriminant $D_{\mathscr{B}/\mathscr{C}}$. Therefore, $\Phi_{\mathscr{C}}^{\circ}\times(\mathbb{C}_{\mathscr{B}\backslash\mathscr{C}}\backslash(Y_{a}\cup Y_{b}))\subset\mathbb{C}_{\mathscr{C}}\circ D_{\mathscr{B}/\mathscr{C}}$
for generic points $\Phi_{\mathscr{C}}^{\circ}\in D_{\mathscr{C}},$
and there exists an inclusion of the closures $D_{\mathscr{C}}\bullet\mathbb{C}_{\mathscr{B}/\mathscr{C}}\subset\mathbb{C}_{\mathscr{C}}\bullet D_{\mathscr{B}/\mathscr{C}}.$
\end{proof}
\begin{thm}
\label{Theorem. Discriminant for a BK-tuple}For a BK-tuple $\mathscr{A}$,
the discriminant $D_{\mathscr{A}}$ is the union $D_{\mathscr{A}}=\underset{\alpha\in\mathtt{P}_{\mathscr{A}}}{\cup}C(\mathscr{B}_{\alpha})$
of strata, enumerated by elements of the poset $\mathtt{P}_{\mathscr{A}}$.
More precisely, the discriminant $D_{\mathscr{A}}$ is a union of
components of codimension one or two: every hypersurface component
corresponds to a stratum $C(\mathscr{B}_{\alpha})$ such that the
tuple $\hat{\mathscr{B}}_{\alpha}$ is a nir; every component of codimension
two corresponds to a stratum $C(\mathscr{B}_{\alpha})$ such that
tuples $\hat{\mathscr{B}}_{\beta}$ are lir for all elements $\beta$
from the order filter $[\alpha]$.
\end{thm}

\begin{proof}
By Corollary \ref{Corollary. Unique decomposition of a BK-tuple},
a BK-tuple $\mathscr{A}$ admits a partition $\mathscr{A}=\underset{\alpha\in\mathtt{P}_{\mathscr{A}}}{\sqcup}\mathscr{B}_{\alpha}$
such that every subtuple $\hat{\mathscr{B}}_{\alpha}$ is an irreducible
BK-tuple and every subtuple $\mathscr{B}_{\mathtt{I}}=\underset{\alpha\in\mathtt{I}}{\sqcup}\mathscr{B}_{\alpha}$
is a BK-tuple for an order ideal $\mathtt{I}\subseteq\mathtt{P}_{\mathscr{A}}$. 

By Corollary \ref{Corollary. Subtuple from the decomposition corresponds to irreducible stratum},
every element $\alpha$ of the poset $\mathtt{P}_{\mathscr{A}}$ corresponds
to a stratum $C(\mathscr{B}_{\alpha}).$ By Separation Lemma \ref{Lemma. Separation Lemma}
and Definition \ref{Definition. BK-multiplication} of BK-multiplication,
the codimension of the stratum $C(\mathscr{B}_{\alpha})$ equals the
codimension of the corresponding discriminant $D_{\hat{\mathscr{B}}_{\alpha}}$
in $\mathbb{C}_{\hat{\mathscr{B}}_{\alpha}}$. By Theorem \ref{Theorem. Codimension of discriminants},
the codimension of $D_{\hat{\mathscr{B}}_{\alpha}}$ equals two/one
if the tuple $\hat{\mathscr{B}}_{\alpha}$ is a lir/nir.

If an element $\alpha$ is covered by an element $\beta$ such that
the tuple $\hat{\mathscr{B}}_{\alpha}$ is a lir and the tuple $\hat{\mathscr{B}}_{\beta}$
is a nir, then we have the inclusion of strata $C(\mathscr{B}_{\alpha})\subset C(\mathscr{B}_{\beta})$
by Lemma \ref{Lemma. If the upper filer consists of linear tuples then strata codim2 lies in codim1}.
Otherwise, the strata do not contain one another and form components
by Lemmas \ref{Lemma. Nir subtuples} and \ref{Lemma. Linear don't lie in each other}.
\end{proof}
\begin{rem}
Denote by $\mathtt{Q}_{\mathscr{A}}$ the order ideal in the poset
$\mathtt{P}_{\mathscr{A}}$, generated by elements $\alpha$ with
a nir $\hat{\mathscr{B}}_{\alpha}$, and its BK-subtuple by $\mathscr{N}=\underset{\alpha\in\mathtt{Q}_{\mathscr{A}}}{\sqcup}\mathscr{B}_{\alpha}$.
By Theorem \ref{Theorem. Discriminant for a BK-tuple}, the discriminant
$D_{\mathscr{N}}$ is a union of hypersurfaces, and the discriminant
$D_{\mathscr{A}/\mathscr{N}}$ is a union of components of codimension
two in the space of linear systems $\mathbb{C}_{\mathscr{A}/\mathscr{N}}$.
\end{rem}

Let us review $\mathscr{A}$-discriminants for dependent tuples $\mathscr{A}$.
Recall that in a matroid, a \textit{circuit} is a minimal dependent
subset, and a \textit{cycle} is a union of circuits.
\begin{thm}
\label{Theorem. A/M  is linearly independent}\textup{\cite{pokidkin_combinatorics_2025}}
For a dependent tuple $\mathscr{A}$, the following are equivalent:

a) $\mathscr{M}$ is the minimal by inclusion subtuple of the minimal
defect;

b) $\mathscr{M}$ is the maximal essential subtuple;

c) $\mathscr{M}$ is the maximal cycle of the induced matroid from
the polymatroid on $\mathscr{A}$.

The subtuple $\mathscr{M}$ is unique, and the quotient $\mathscr{A}/\mathscr{M}$
is independent.
\end{thm}

Theorem \ref{Theorem. A/M  is linearly independent} provides a new
characterization of the subtuple $\mathscr{M}$ for the resultant
\cite{sturmfels_newton_1994} as the maximal cycle of the induced
matroid from the realizable polymatroid.
\begin{thm}
For a dependent tuple $\mathscr{A}$ with a maximal essential subtuple
$\mathscr{M}$, the sparse resultant $R_{\mathscr{A}}$ equals the
sparse resultant $R_{\mathscr{M}}$ of codimension $-\delta(\mathscr{M})$.
\end{thm}

\begin{proof}
See algebro-geometric proofs in Theorem 1.1 \cite{sturmfels_newton_1994}
and Theorem 2.15 \cite{esterov_determinantal_2007}, and a tropical
proof in Theorem 2.23 \cite{jensen_computing_2013}.
\end{proof}
\begin{cor}
\label{Corollary. Discriminants for linearly dependent tuples}For
a dependent tuple $\mathscr{A}$ with a maximal essential subtuple
$\mathscr{M}$, the discriminant $D_{\mathscr{A}}$ is the sparse
resultant $R_{\mathscr{M}}$.
\end{cor}


\section{\label{Section. Cayley discriminants}Cayley discriminants}

The \textit{Cayley discriminant} $D_{cay(\mathscr{A})}$ is the discriminant
for the Cayley set $cay(\mathscr{A})$. If a BK-tuple $\mathscr{A}$
is irreducible, then the $\mathscr{A}$-discriminant equals the Cayley
discriminant (Corollary \ref{Corollary. Irreducible BK =00003D> three types of discriminant are equal}).

Notice the isomorphism of vector spaces $\mathbb{C}_{\mathscr{A}}\cong\mathbb{C}_{cay(\mathscr{A})}.$
\begin{lem}
\label{Lemma. Cayley discriminant for a connected component of the poset}If
the poset $\mathtt{P}_{\mathscr{B}}$ is a connected component of
the poset $\mathtt{P}_{\mathscr{A}}$ for BK-tuples $\mathscr{B}\subset\mathscr{A}$,
then the Cayley discriminant is the intersection
\[
D_{cay(\mathscr{A})}\cong D_{cay(\mathscr{B})}\bullet\mathbb{C}_{\mathscr{A}/\mathscr{B}}\cap\mathbb{C}_{\mathscr{B}}\bullet D_{cay(\mathscr{A}/\mathscr{B})}.
\]
\end{lem}

\begin{proof}
Since the poset $\mathtt{P}_{\mathscr{B}}$ is a connected component,
the tuple $\mathscr{A}\backslash\mathscr{B}$ is a BK-tuple. Hence
the linear span of the BK-tuple $\mathscr{A}$ decomposes into the
direct sum $\langle\mathscr{A}\rangle=\langle\mathscr{B}\rangle\oplus\langle\mathscr{A}\backslash\mathscr{B}\rangle$,
and the torus splits $T(\mathscr{A})=T(\mathscr{B})\times T(\mathscr{A}\backslash\mathscr{B}).$
Then every system $\Phi_{\mathscr{A}}(x,y)$ from $\mathbb{C}_{\mathscr{A}}$
splits into two independent systems $\Phi_{\mathscr{B}}(x)$ and $\Phi_{\mathscr{A}\backslash\mathscr{B}}(y)$,
and every polynomial $f\in\mathbb{C}_{cay(\mathscr{A})}$ can be written
\[
f(x,y,\lambda)=\underset{A\in\mathscr{A}}{\sum}\lambda_{A}f_{A}(x,y)=\underset{B\in\mathscr{B}}{\sum}\lambda_{B}f_{B}(x)+\underset{C\in\mathscr{A}\backslash\mathscr{B}}{\sum}\lambda_{C}f_{C}(y).
\]
This means that a point $(x^{\circ},y^{\circ},\lambda^{\circ})$ is
a singular point for the polynomial $f_{cay(\mathscr{A})}$ if and
only if $(x^{\circ},\lambda_{\mathscr{B}}^{\circ})$ and $(y^{\circ},\lambda_{\mathscr{\mathscr{A}\backslash B}}^{\circ})$
are singular points for polynomials $f_{cay(\mathscr{B})}$ and $f_{cay(\mathscr{A}\backslash\mathscr{B})}$. 
\end{proof}
\begin{lem}
\label{Lemma. Simplification of a Cayley discriminant}For a BK-tuple
$\mathscr{A}$ with a connected poset $\mathtt{P}_{\mathscr{A}}$
and a BK-subtuple $\mathscr{B}$, the Cayley discriminant is isomorphic
to $D_{cay(\mathscr{A})}\cong\mathbb{C}_{\mathscr{B}}\bullet D_{cay(\mathscr{A}/\mathscr{B})}$.
\end{lem}

\begin{proof}
We can choose a basis such that systems from $\mathbb{C}_{\mathscr{B}}$
depend on variables $x\in T(\mathscr{B})$ and systems from $\mathbb{C}_{\mathscr{A}\backslash\mathscr{B}}$
depend on variables $(x,y)\in T(\mathscr{A}).$ Then every polynomial
$f\in\mathbb{C}_{cay(\mathscr{A})}$ has the form: $f(x,y,\lambda)=\underset{B\in\mathscr{B}}{\sum}\lambda_{B}f_{B}(x)+\underset{C\in\mathscr{A}\backslash\mathscr{B}}{\sum}\lambda_{C}f_{C}(x,y).$

$\boxed{\subseteq}$ A generic polynomial $f(x,y,\lambda)$ from the
discriminant $D_{cay(\mathscr{A})}$ has a singular point $(x^{\circ},y^{\circ},\lambda^{\circ})$
in the torus $T(cay(\mathscr{A}))$ such that
\begin{align*}
\nabla_{y}f(x^{\circ},y^{\circ},\lambda^{\circ}) & =\underset{C\in\mathscr{A}\backslash\mathscr{B}}{\sum}\lambda_{C}^{\circ}\,\nabla_{y}f_{C}(x^{\circ},y^{\circ})=\underset{C\in\mathscr{A}\backslash\mathscr{B}}{\sum}\lambda_{C}^{\circ}\,\nabla_{y}(ev_{x^{\circ}}f_{C})(y^{\circ})=0.
\end{align*}
This means the point $(y^{\circ},\lambda_{\mathscr{A}\backslash\mathscr{B}}^{\circ})$
is singular for the polynomial $f_{cay(\mathscr{A}/\mathscr{B})}=\underset{B\in\mathscr{A}\backslash\mathscr{B}}{\sum}\lambda_{B}^{\circ}\,(ev_{x^{\circ}}f_{B})(y)$,
and the corresponding system $(\Phi_{\mathscr{B}},f_{cay(\mathscr{A}\backslash\mathscr{B})})$
belongs to $\mathbb{C}_{\mathscr{B}}\bullet D_{cay(\mathscr{A}/\mathscr{B})}$.

$\boxed{\supseteq}$ For a generic system $(\Phi_{\mathscr{B}},f_{cay(\mathscr{A}\backslash\mathscr{B})})$
from $\mathbb{C}_{\mathscr{B}}\circ D_{cay(\mathscr{A}/\mathscr{B})}$,
there exists a root $(x^{\circ},y^{\circ},\lambda_{\mathscr{A}\backslash\mathscr{B}}^{\circ})$
from the torus $T(\mathscr{B})\times T(cay(\mathscr{A}/\mathscr{B}))$
such that $\underset{C\in\mathscr{A}\backslash\mathscr{B}}{\sum}\lambda_{C}^{\circ}\,\nabla_{y}(ev_{x^{\circ}}f_{C})(y^{\circ})=0$. 

Since the poset $\mathtt{P}_{\mathscr{A}}$ is connected, the intersection
of saturated sublattices $\mathscr{\langle\mathscr{B}}\rangle$ and
$\langle\mathscr{A}\backslash\mathscr{B}\rangle$ has a positive dimension.
Hence the vector $\underset{C\in\mathscr{A}\backslash\mathscr{B}}{\sum}\lambda_{C}^{\circ}\,\nabla_{x}(ev_{y^{\circ}}f_{C})(x^{\circ})$
is not zero, and there exists a nontrivial linear combination of linearly
dependent vectors
\[
\underset{B\in\mathscr{B}}{\sum}\mu_{B}\,\nabla_{x}f_{B}(x^{\circ})+\mu\underset{C\in\mathscr{A}\backslash\mathscr{B}}{\sum}\lambda_{C}^{\circ}\,\nabla_{x}(ev_{y^{\circ}}f_{C})(x^{\circ})=0,
\]
for some constants $\{\mu_{B}\}_{B\in\mathscr{B}}$ and $\mu$. These
constants are non-zero because vectors $\{\nabla_{x}f_{B}(x^{\circ})\}_{B\in\mathscr{B}}$
are linearly independent for the generic system $\Phi_{\mathscr{B}}$.
Therefore, the corresponding Cayley polynomial $f_{cay(\mathscr{A})}$
has the singular point $(x^{\circ},y^{\circ},\mu_{\mathscr{B}},\mu\lambda_{\mathscr{A}\backslash\mathscr{B}}^{\circ})$
in the torus $T(cay(\mathscr{A}))$ and belongs to the discriminant
$D_{cay(\mathscr{A})}$.
\end{proof}
\begin{prop}
\label{Proposition. Discriminant's strata are Cayley discriminants}For
an element $\alpha$ of the poset $\mathtt{P}_{\mathscr{A}}$, the
stratum $C(\mathscr{B}_{\alpha})$ of the $\mathscr{A}$-discriminant
can be expressed via the Cayley discriminant: $D_{cay(\mathscr{B}_{(\alpha)})}=\mathbb{C}_{\mathscr{B}_{(\alpha)\backslash\alpha}}\bullet D_{\hat{\mathscr{B}}_{\alpha}}.$
\end{prop}

\begin{proof}
The order ideal $(\alpha)$ has the unique maximal element $\alpha.$
Iteratively apply Lemma \ref{Lemma. Simplification of a Cayley discriminant}
and use the coherency relations \ref{Corollary. Coherency Relations}.
At the last step, we obtain the Cayley discriminant of the irreducible
BK-tuple $\hat{\mathscr{B}}_{\alpha},$ which coincides with the $\mathscr{A}$-discriminant
by Theorem \ref{Theorem. Codimension of discriminants}.
\end{proof}
\begin{rem}
This proposition sets up the consistency of Theorem \ref{Theorem. Discriminant for a BK-tuple}
about components of the $\mathscr{A}$-discriminant for a BK-tuple
with Esterov Theorem 2.31 \cite{esterov_newton_2010} about components
of the reduced discriminant, expressed as Cayley discriminants.
\end{rem}

\begin{thm}
\label{Theorem. Cayley Discriminants}For a BK-tuple $\mathscr{A}$,
the Cayley discriminant $D_{cay(\mathscr{A})}$ equals the intersection
of $\mathscr{A}$-discriminant components, enumerated by maximal elements
of the poset $\mathtt{P}_{\mathscr{A}}$
\begin{align*}
D_{cay(\mathscr{A})} & =\underset{\alpha\in max\,\mathtt{P}_{\mathscr{A}}}{\cap}\mathbb{C}_{\mathscr{\mathscr{A}}\backslash\mathscr{B}_{\alpha}}\bullet D_{\hat{\mathscr{B}}_{\alpha}}.
\end{align*}
The intersection is complete and of codimension $2l+n$, where $l$/$n$
is the number of all maximal elements $\alpha$ in the poset $\mathtt{P}_{\mathscr{A}}$
such that the BK-tuple $\hat{\mathscr{B}}_{\alpha}$ is a lir/nir.
\end{thm}

\begin{proof}
By Lemma \ref{Lemma. Simplification of a Cayley discriminant} and
coherency relations from Corollary \ref{Corollary. Coherency Relations},
we can simplify the Cayley discriminant for the BK-tuple $\mathscr{A}$
to the Cayley discriminant for the BK-tuple $\mathscr{A}/\mathscr{C}$,
where $\mathscr{C}=\mathscr{A}\backslash\underset{\alpha\in max\,\mathtt{P}_{\mathscr{A}}}{\cup}\mathscr{B}_{\alpha}$.
Then the poset $\mathtt{P}_{\mathscr{A}/\mathscr{C}}$ is a disjoint
union of $|max\,\mathtt{P}_{\mathscr{A}}|$ incomparable elements.
This means that the BK-tuple $\mathscr{A}/\mathscr{C}$ is semi-irreducible,
and $\mathbb{C}_{\mathscr{A}/\mathscr{C}}=\underset{\alpha\in max\,\mathtt{P}_{\mathscr{A}}}{\prod}\mathbb{C}_{\hat{\mathscr{B}}_{\alpha}}$.
Hence the Cayley discriminant $D_{cay(\mathscr{A}/\mathscr{C})}$
equals the direct product $\underset{\alpha\in max\,\mathtt{P}_{\mathscr{A}}}{\prod}D_{\hat{\mathscr{B}}_{\alpha}}$
by Lemma \ref{Lemma. Cayley discriminant for a connected component of the poset}
and by using the isomorphism of discriminants $D_{cay(\hat{\mathscr{B}}_{\alpha})}=D_{\hat{\mathscr{B}}_{\alpha}}$
for irreducible BK-tuples $\hat{\mathscr{B}}_{\alpha}$. Therefore,
the intersection $\underset{\alpha\in max\,\mathtt{P}_{\mathscr{A}}}{\cap}\mathbb{C}_{\mathscr{\mathscr{A}}\backslash\mathscr{B}_{\alpha}}\bullet D_{\hat{\mathscr{B}}_{\alpha}}$
is complete, and the codimension is clear.
\end{proof}
Following the proof of Lemma \ref{Lemma. Simplification of a Cayley discriminant},
it is possible to show (notice Theorem \ref{Theorem. A/M  is linearly independent}):
\begin{thm}
\label{Theorem. The Cayley discriminant for a linearly dependent tuple}For
a dependent tuple $\mathscr{A},$ the Cayley discriminant equals the
multiplication $R_{\mathscr{M}}\bullet D_{cay(\mathscr{A}/\mathscr{M})}$,
where $R_{\mathscr{M}}$ is the resultant of the maximal essential
subtuple $\mathscr{M}$.
\end{thm}

\begin{rem}
The subtuple $\mathscr{A}\backslash\mathscr{M}$ consists of coloops
of the induced matroid from the polymatroid on $\mathscr{A}$, and
the quotient $\mathscr{A}/\mathscr{M}$ is independent by Theorem
\ref{Theorem. A/M  is linearly independent}. Since additional coefficients
in the Cayley trick are not allowed to be zero, this clarifies the
distinguished contribution of $D_{cay(\mathscr{A}/\mathscr{M})}$
in the BK-multiplication. 
\end{rem}

\begin{cor}
\label{Corollary. Cayley discriminant for an essential linearly dependent tuple}For
an essential dependent tuple $\mathscr{A}$, the Cayley discriminant
equals the sparse resultant $R_{\mathscr{A}}$ \textup{(see Proposition
6.1 and Lemma 6.3 \cite{dickenstein_tropical_2007})}.
\end{cor}


\section{\label{Section. Mixed discriminants}Mixed discriminants}
\begin{lem}
\label{Lemma. Mixed discriminant for a connected component of the poset}If
the poset $\mathtt{P}_{\mathscr{A}}$ has more than one connected
component for a BK-tuple $\mathscr{A}$, then the mixed discriminant
is empty.
\end{lem}

\begin{proof}
Every connected component corresponds to a BK-subtuple of a BK-tuple
$\mathscr{A}$, and the tuple $\mathscr{A}$ decomposes into a disjoint
union of BK-subtuples. Then we can split the variables and polynomial
systems into independent subsystems with BK-supports. This means that
a polynomial system with the support $\mathscr{A}$ cannot have a
non-degenerate multiple root.
\end{proof}
\begin{lem}
\label{Lemma. Simplification of a mixed discriminant}For BK-tuples
$\mathscr{B}\subset\mathscr{A}$, the mixed discriminant $\mathring{D}_{\mathscr{A}}$
equals $\mathbb{C}_{\mathscr{B}}\bullet\mathring{D}_{\mathscr{A}/\mathscr{B}}$.
\end{lem}

\begin{proof}
$\boxed{\subseteq}$ If the mixed discriminant $\mathring{D}_{\mathscr{A}}$
is not empty, then its generic system $\Phi$ has a non-degenerate
multiple root $(x,y)\in T(\mathscr{A})$ such that vectors $\{df_{A}(x,y)\}_{A\in\mathscr{A}}$
are linearly dependent, and $x\in T(\mathscr{B})$. If the point $x$
is singular for the subsystem $\{df_{B}(x)\}_{B\in\mathscr{B}}$,
then the point $(x,y)$ cannot be a non-degenerate multiple root for
the system $\Phi$. Hence the set $\{d(ev_{x}f_{C})(y)\}_{C\in\mathscr{A}/\mathscr{B}}$
is linearly dependent by Lemma \ref{Lemma. Linear Algebra}. Moreover,
the point $y$ is a non-degenerate multiple root for the system $ev_{x}\Phi_{\mathscr{A}\backslash\mathscr{B}}$,
and we obtain the inclusion for discriminants. $\boxed{\supseteq}$
It is clear. 
\end{proof}
\begin{thm}
\label{Theorem. The mixed discriminant for a BK-tuple}For a BK-tuple
$\mathscr{A}$, the mixed discriminant equals the Cayley discriminant
if the poset $\mathtt{P}_{\mathscr{A}}$ has only one maximal element.
Otherwise, the mixed discriminant is empty.
\end{thm}

\begin{proof}
By Lemma \ref{Lemma. Simplification of a mixed discriminant} and
coherency relations from Corollary \ref{Corollary. Coherency Relations},
we can simplify the mixed discriminant for the BK-tuple $\mathscr{A}$
to the mixed discriminant for the BK-tuple $\mathscr{A}/\mathscr{C}$,
where $\mathscr{C}=\mathscr{A}\backslash\underset{\alpha\in max\,\mathtt{P}_{\mathscr{A}}}{\cup}\mathscr{B}_{\alpha}$.
Then the poset $\mathtt{P}_{\mathscr{A}/\mathscr{C}}$ is a disjoint
union of $|max\,\mathtt{P}_{\mathscr{A}}|$ incomparable elements.
The mixed discriminant $\mathring{D}_{\mathscr{A}/\mathscr{C}}$ is
not empty only if the poset $\mathtt{P}_{\mathscr{A}/\mathscr{C}}$
consists of one element $\alpha$ by Lemma \ref{Lemma. Mixed discriminant for a connected component of the poset}.
In that case, the BK-tuple $\mathscr{A}/\mathscr{C}$ is irreducible,
and the three types of discriminants coincide. Therefore, using Proposition
\ref{Proposition. Discriminant's strata are Cayley discriminants}
for the BK-tuple $\mathscr{A}$, the mixed discriminant and the Cayley
discriminant are equal to $\mathbb{C}_{\mathscr{C}}\bullet D_{\mathscr{A}/\mathscr{C}}$.
\end{proof}
\begin{rem}
The paper \cite{cattani_mixed_2013} shows that if the Cayley discriminant
for a BK-tuple is a hypersurface, then the mixed discriminant is the
same hypersurface. Theorem \ref{Theorem. The mixed discriminant for a BK-tuple}/Theorem
\ref{Theorem. The mixed discriminant for a linearly dependent tuple}
provides the second proof of this result for nonlinear BK-tuples/dependent
tuples.
\end{rem}

\begin{thm}
\label{Theorem. The mixed discriminant for a linearly dependent tuple}The
mixed discriminant of a dependent tuple $\mathscr{A}$ is not empty
if and only if the tuple contains only one circuit $\mathscr{C}$.
Then the mixed discriminant is the resultant $R_{\mathscr{C}}$, and
it is a hypersurface.
\end{thm}

\begin{proof}
If a tuple $\mathscr{A}$ contains more than one circuit, then the
tuples $\mathscr{A}\backslash A$ are dependent for each set $A\in\mathscr{A}$.
This means that systems from the space $\mathbb{C}_{\mathscr{A}}$
cannot have a non-degenerate multiple root, and the mixed discriminant
is empty.

If the tuple $\mathscr{A}$ contains only one circuit $\mathscr{C}$,
then every tuple $\mathscr{A}\backslash C$ is independent for every
set $C\in\mathscr{C}$. Then every root for a polynomial system from
$\mathbb{C}_{\mathscr{A}}$ is a non-degenerate multiple root, and
the mixed discriminant coincides with the resultant $R_{\mathscr{C}}$.
Since circuits in the induced matroid always have the defect $-1$
\cite{pokidkin_combinatorics_2025}, the mixed discriminant is a hypersurface.
\end{proof}

\section{\label{Section. Degrees of components}Degrees of discriminants}
\begin{cor}
\textup{\label{Corollary. Degree of the mixed discriminant for a linearly dependent tuple}\cite{pedersen_product_1993}
}For a dependent tuple with the unique circuit $\mathscr{C}$, the
mixed discriminant has the degree $\underset{C\in\mathscr{C}}{\sum}\mathrm{MV}_{\overline{\langle\mathscr{C}\rangle}}(\mathscr{C}\backslash C)$.
\end{cor}

For an essential dependent tuple $\mathscr{A}$, the $\mathscr{A}$-discriminant
and the Cayley discriminant are equal to the sparse resultant $R_{\mathscr{A}}$
by Corollary \ref{Corollary. Discriminants for linearly dependent tuples}
and Corollary \ref{Corollary. Cayley discriminant for an essential linearly dependent tuple}.
We present two formulas to compute the degree of a sparse resultant.
Notice that the sparse resultant of an independent tuple equals the
sparse resultant of its maximal essential subtuple.
\begin{prop}
\textup{\label{Proposition. I Degree of the sparse resultant for a linearly dependent tuple}}For
an essential dependent tuple $\mathscr{A}$, the sparse resultant
$R_{\mathscr{A}}$ has the degree $\underset{\mathscr{B}\in\mathcal{B}(\mathscr{A})}{\sum}\mathrm{MV}_{\overline{\langle\mathscr{A}\rangle}}(\mathscr{B}),$
where $\mathcal{B}(\mathscr{A})$ is the set of bases of the induced
matroid on $\mathscr{A}$.
\end{prop}

\begin{proof}
By Corollary 6.5 in \cite{dickenstein_tropical_2007}, the degree
of the sparse resultant $R_{\mathscr{A}}$ is the sum of mixed volumes
$\mathrm{MV}_{\overline{\langle\mathscr{A}\rangle}}(\mathscr{B})$
over all subtuples $\mathscr{B}$ of cardinality $dim\,\langle\mathscr{A}\rangle$.
For the induced matroid on $\mathscr{A}$, the bases are BK-tuples
of cardinality $dim\,\langle\mathscr{A}\rangle$ according to \cite{pokidkin_combinatorics_2025}.
Hence all subtuples of cardinality $dim\,\langle\mathscr{A}\rangle$
are either dependent or bases and BK-tuples. Therefore, in the formula
from \cite{dickenstein_tropical_2007}, only bases of the induced
matroid contribute to the degree of the sparse resultant $R_{\mathscr{A}}$.
\end{proof}
Let $\mathscr{A}$ be an essential dependent tuple of defect $-\delta$.
Construct the lattice $L=\overline{\langle\mathscr{A}\rangle}\times\mathbb{Z}^{\delta}$
and the new tuple $\mathscr{A}^{\flat}=(A\times\triangle_{\delta},\,A\in\mathscr{A})$,
where $\triangle_{\delta}$ is the standard simplex in $\mathbb{Z}^{\delta}$,
and each set $A\times\triangle_{\delta}$ is considered in the lattice
$L$. Notice that the tuple $\mathscr{A}^{\flat}$ has zero defect.
\begin{prop}
\textup{\label{Proposition. II Degree of the sparse resultant for a linearly dependent tuple}(Esterov)}
For an essential dependent tuple $\mathscr{A}$, the sparse resultant
$R_{\mathscr{A}}$ has the degree $\mathrm{MV}_{L}(\mathscr{A}^{\flat})$.
\end{prop}

\begin{proof}
The degree of the resultant $R_{\mathscr{A}}$ with codimension $\delta$
equals the number of intersection points with a generic $\delta$-dimensional
vector subspace $\Pi$ in $\mathbb{C}_{\mathscr{A}}$. Choose a parametrization
for the subspace $\Pi$: $\Phi=\text{\ensuremath{\Phi}}^{0}+\stackrel[i=1]{\delta}{\sum}y_{i}\Phi^{i}$
for some fixed choice of points $\Phi^{0},...,\Phi^{\delta}$ from
$\mathbb{C}_{\mathscr{A}}$ and new variables $y_{1},...,y_{\delta}$.
Each coefficient of the system $\Phi$ is a linear function of the
new variables $y$, and the parametrization provides a polynomial
system from the support $\mathbb{C}_{\mathscr{A}^{\flat}}$. By the
Kouchnirenko-Bernstein theorem, a generic system from $\mathbb{C}_{\mathscr{A}^{\flat}}$
has $\mathrm{MV}_{L}(\mathscr{A}^{\flat})$ solutions. Each solution
corresponds to an intersection of the hyperplane $\Pi$ with the resultant.

Since the tuple $\mathscr{A}$ is essential, the tuple $\mathscr{A}^{\flat}$
is irreducible because every proper subtuple $\mathscr{B}^{\flat}$
has a positive defect: $\delta(\mathscr{B}^{\flat})=\delta(\mathscr{B})-\delta(\mathscr{A})>0$.
This observation ensures that the mixed volume $\mathrm{MV}_{L}(\mathscr{A}^{\flat})$
is always positive.
\end{proof}
For a face $A'$ of a set $A$ from a lattice$,$ consider the projection
of saturated sublattices $s:\ensuremath{\overline{\langle A\rangle}}\rightarrow\ensuremath{\overline{\langle A\rangle}}/\text{\ensuremath{\overline{\langle A'\rangle}}}$
and the numbers $c^{A',A}$, which are differences of integer volumes
between sets $s(A)$ and $s(A\backslash A')$ in the lattice $s(\ensuremath{\overline{\langle A\rangle}}).$
Set $c^{A,A}=1$ and $c^{A',A}=0$ if $A'$ is not a face of $A$.
Then we can define a square matrix $C$ with entries $c^{A'',A'}$
for all possible faces $A''$, $A'$, and build the inverse matrix
$C^{-1}$ with entries $e^{A'',A'}$, called \textit{Euler obstructions}
(see \cite{esterov_newton_2010}).
\begin{thm}
\label{Theorem. Matsui-Takeuchi, Degree of the Cayley discriminant}\textup{\cite{matsui_geometric_2011}
For a finite set $A$, the $A$-discriminant of codimension $\delta$
has the degree
\[
deg\,D_{A}=\underset{A'\subseteq A}{\sum}e^{A',A}\,\left(\binom{dim\,\langle A'\rangle-1}{\delta}+(-1)^{\delta+1}(\delta+1)\right)\,\mathrm{Vol}(A').
\]
}
\end{thm}

For a codimension one $A$-discriminant, the formula coincides with
results \cite{gelfand_discriminants_1994,dickenstein_tropical_2007}.
In particular, this theorem describes degrees of Cayley discriminants
for the Cayley set $cay(\mathscr{A})$ of a tuple $\mathscr{A}.$
If the tuple $\mathscr{A}$ is dependent with a maximal essential
subtuple $\mathscr{M}$, then we use the formula for the codimension
$\delta=-\delta(\mathscr{M})$. If the tuple $\mathscr{A}$ is BK,
then we use the formula for the codimension $\delta=2l+n$ by Theorem
\ref{Theorem. Cayley Discriminants}. Moreover, for a BK-tuple $\mathscr{A}$,
we can use Theorem \ref{Theorem. Matsui-Takeuchi, Degree of the Cayley discriminant}
to compute the degrees for components of the $\mathscr{A}$-discriminant. 
\begin{cor}
\label{Corollary. Degrees for components of the discriminant for a BK-tuple}For
a BK-tuple $\mathscr{A}$ and $\alpha\in P_{\mathscr{A}},$ the component
$C(\mathscr{B}_{\alpha})$ has the degree 
\begin{align*}
deg\,C(\mathscr{B}_{\alpha}) & =\begin{cases}
\underset{A\subseteq cay(\mathscr{B}_{(\alpha)})}{\sum}e^{A,cay(\mathscr{B}_{(\alpha)})}\,(dim\,\langle A\rangle+1)\,\mathrm{Vol}(A), & \text{if }\hat{\mathscr{B}}_{\alpha}\text{ is a nir,}\\
\frac{1}{2}\underset{A\subseteq cay(\mathscr{B}_{(\alpha)})}{\sum}e^{A,cay(\mathscr{B}_{(\alpha)})}\,(dim\,\langle A\rangle+1)\,(dim\,\langle A\rangle-4)\,\mathrm{Vol}(A), & \text{if }\hat{\mathscr{B}}_{\alpha}\text{ is a lir.}
\end{cases}
\end{align*}
\end{cor}

\begin{rem}
\label{Remark. Degree of the mixed discriminant for a BK-tuple}1)
Corollary \ref{Corollary. Degrees for components of the discriminant for a BK-tuple}
describes degrees of mixed discriminants for BK-tuples $\mathscr{A}$
such that the poset $\mathtt{P}_{\mathscr{A}}$ has a unique maximal
element $\alpha$, $\mathscr{A}=\mathscr{B}_{(\alpha)}$.

2) For a nir $\hat{\mathscr{B}}_{\alpha}$ and its face $A$, the
number $(dim\,\langle A\rangle+1)\,\mathrm{Vol}(A)$ is the total
degree of the $A$-Euler discriminant (see Corollary 1.9 \cite{esterov_discriminant_2013}).
\end{rem}

Every face of a Cayley set is a Cayley set for some collection of
faces. Esterov expressed volumes of Cayley sets (see Lemma 1.7 \cite{esterov_newton_2010})
and mixed volumes for tuples of Cayley sets (see \cite{esterov_multiplicities_2012})
via mixed volumes of their generating sets. The computation of volumes
simplifies the Matsui-Takeuchi degree formula (see Corollary 1.14
\cite{esterov_discriminant_2013}).

Denote the mixed volume of $n$ finite sets by a monomial $A_{1}\cdot...\cdot A_{n}$
and the integer simplex $\{a\in\mathbb{Z}_{\geq0}^{k}\,|\;a_{1}+...+a_{k}=m\}$
by $\triangle_{k}(m)$.
\begin{cor}
\textup{\cite{esterov_newton_2010}} For a tuple of finite sets $\mathscr{A}=(A_{1},...,A_{k})$,
the volume of the Cayley set $cay(\mathscr{A})$ in its $n$-dimensional
linear span equals
\[
\mathrm{Vol}(cay(\mathscr{A}))=\underset{a\in\triangle_{k}(n)}{\sum}A_{1}^{a_{1}}\cdot...\cdot A_{k}^{a_{k}}.
\]
\end{cor}

\begin{cor}
\label{Corollary. Degree of the Cayley discriminant for a BK-tuple}For
a BK-tuple $\mathscr{A}$, the Cayley discriminant has the degree
$\underset{\alpha\in max\,\mathtt{P}_{\mathscr{A}}}{\prod}deg\,C(\mathscr{B}_{\alpha})$.
\end{cor}

\begin{proof}
We have a complete intersection by Theorem \ref{Theorem. Cayley Discriminants}.
\end{proof}
\begin{prop}
For a lir $\mathscr{A}$ of cardinality $n$, the $\mathscr{A}$-discriminant
has the degree $\frac{n\,(n+1)}{2}$.
\end{prop}

\begin{proof}
For a linear BK-tuple $\mathscr{A},$ the discriminant $D_{\mathscr{A}}$
is a determinantal variety. One of the proofs is written in Example
19.10 \cite{harris_algebraic_1992}.
\end{proof}
\begin{problem*}
The description of components and degrees of mixed and $\mathscr{A}$-discriminants
is still open for underdetermined polynomial systems consisting of
more than one equation.
\end{problem*}

\subsubsection*{Acknowledgment}

I am deeply grateful to Alexander Esterov for supervision and to my
family for support.

\subsubsection*{Funding}

The study was partially supported by the HSE University Basic Research
Program and by the Foundation for the Advancement of Theoretical Physics
and Mathematics \textquotedblleft BASIS\textquotedblright .

\begin{adjustwidth}{-0.0in}{-0.0in}
	\sloppy
	\printbibliography
	\fussy
\end{adjustwidth}


\end{document}
